%\subsection{Caracterización del microbioma}
%% Que hay en el microbioma?--> Taxones--> Focus en bacterias %%
%%% Explicar que hay en el microbioma, como los identificamos, hemos centrado la tesis solo en bacterias

Todos los organismos vivos, incluidos los microorganismos, se organizan según su similitud o relación evolutiva en diferentes categorías taxonómicas. Estas categorías, de mayor a menor, son: dominio, filo, clase, orden, familia, género y especie. Cada nivel está contenido dentro del anterior. Los dominios incluyen grupos de filos, las clases agrupan órdenes, los órdenes incluyen familias, las familias agrupan géneros, y los géneros contienen especies \cite{pitt2012classification}. La especie es la categoría más específica y se define como un grupo de organismos estrechamente relacionados que pueden reproducirse entre sí en eucariotas. En procariotas, se define como una población de células con características similares.

Las bacterias, en particular, son componentes fundamentales del microbioma humano y poseen una variedad de formas y funciones que hacen posible su adaptación a diferentes entornos corporales, como la piel, la boca, el tracto gastrointestinal, entre otros \cite{cho2012human}. Dentro del microbioma humano, las bacterias participan en una multitud de procesos vitales. Ayudan en la digestión de compuestos complejos, la síntesis de vitaminas esenciales, y la modulación del sistema inmune. Además, actúan como una barrera defensiva, inhibiendo el crecimiento de patógenos mediante la competencia por recursos y la producción de sustancias antimicrobianas \cite{backhed2005host}.

En muchos casos, ya a nivel de filo (segundo nivel taxonomico) se puede observar diferencias signifcativas en la composición microbiana. Por ejemplo, en el caso del intestino humano los filos más predominantes son Firmicutes y Bacteroidetes (juntos pueden constituir más del 90\% de la comunidad bacteriana intestinal) \cite{ley2006human}. Estas diferencias en la composición a nivel de filo no solo impactan las funciones metabólicas, sino también el equilibrio general del microbioma. Un equilibrio adecuado entre Firmicutes y Bacteroidetes (también en menor medido el resto de filos) es esencial para mantener la salud del huésped, fomentando la digestión eficiente y la protección inmunológica \cite{turnbaugh2009core}.



%%%% Introducir el termino de metagenómica. Hablar de los tipos de muestras de las que podemos extraer una comunidad microbiana. Pinceladas pasos y protocolos. NGS y tipos de secuenciación. WGS vs Amplicon 16S %%%

Para el estudio del microbioma humano acudimos a la secuenciación metagenómica, esto es, el estudio del material genético de un ambiente microbiano sin necesidad de realizar un cultivo (aunque en el pasado este término hacia alusión solamente a la metodología de secuenciación metagenómica de escopeta).

Un paso crucial de cualquier estudio para la búsqueda de biomarcadores es la selección del tipo de muestra a utilizar. Esta elección es crucial y debe alinearse con los objetivos específicos del estudio y la naturaleza de la condición a investigar. Las muestras microbianas pueden variar ampliamente según la fuente de la infección o la necesidad del análisis. Las muestras más comunes incluyen sangre, orina, heces, esputo, frotis y tejidos. Cada tipo de muestra ofrece ventajas y desventajas únicas en términos de accesibilidad, manejabilidad, y la riqueza de información que pueden proporcionar sobre los microorganismos presentes y su potencial como biomarcadores \cite{claesson2017clinician}. Otros aspectos a tener en cuenta son la extracción, manipulación, almacenamiento y extracción del ADN de la muestra microbiana. La manipulación de estas nuestras tras su recogida es un aspecto crucial en los métodos de secuenciación basados en ácidos nucleicos para comparar la composición y diversidad de las comunidades microbianas. De hecho, la recogida de muestras podría ser un importante factor de confusión que podría afectar a los resultados y las interpretaciones del estudio \cite{thomas2012metagenomics}. Para la correcta realización de estos pasos existe un amplio catálogo de protcolos en el estado del arte, así como, un gran número de kits comerciales para la realización de la extracción del ADN \cite{claesson2017clinician}.



\textbf{Unidades taxonomicas operativas}

Las OTUs son agrupaciones de secuencias genéticas que se consideran equivalentes, o lo suficientemente similares, para clasificarlas bajo una misma 'especie' en estudios de microbiomas. Generalmente, estas agrupaciones se forman en función de un umbral de similitud, típicamente alrededor del 97\% para especie y 95\% para género. Esto implica que todas las secuencias dentro de una OTU tienen al menos esa proporción de similitud. Facilitan la comparación de distintos conjuntos de datos y simplifican la complejidad del análisis del microbioma. Dado que se basan en un umbral arbitrario de similitud, pueden incluir variabilidad no biológica, lo que a veces lleva a perder detalles importantes sobre la diversidad \cite{schloss2011assessing, edgar2018updating}.

\textbf{Variantes de secuencia de amplicón}

Las ASVs representan secuencias únicas exactas de la región del amplicón, sin agruparlas en función de la similitud. Esto significa que cada ASV es una secuencia distinta, incluso si difiere por un solo nucleótido. Ofrecen una resolución mucho más alta que las OTUs, permitiendo identificar cambios sutiles en la composición del microbioma. Permiten comparaciones más detalladas y precisas entre diferentes estudios y bases de datos, ya que no están limitadas por umbrales de similitud. Requieren métodos más avanzados de procesamiento y análisis de datos para garantizar la precisión y evitar errores de secuenciación o artefactos \cite{callahan2016dada2, callahan2017exact}


\subsection{Generación y almacenamiento de datos}

Quiero hablar del HMP, ya que fue el punto de partida de la investigacion del microbioma humano, sentando las bases de un microbioma "sano" baselina para estudiar los procesos disbioticos que pueden estar interfiriendo en enfermedades.  Además hablar de que gracias al FAIR, muchos autores están compartiendo tanto los datos crudos como las matrices de conteos para que la comunidad científica pueda reutilizar esos datos en sus investigaciones. Por ende, surje la necesidad de que existan bases de datos curated e introducir las que considero más importantes y las mas utilizadas: GenBank, ENA y DDBJ. ADemás este triangulo de bases de datos conforma una colaboracion internacionales para mantener, actualizar y estandarizar los datos para toda la comunidad cientifica global.

\subsubsection{Iniciativas}

Como ya hemos comentado previamente el HMP y su fase integrativa posterior representan iniciativas fundamentales que han avanzado en la investigación del microbioma humano mediante la generación, integración y compartición abierta de grandes conjuntos de datos multi-ómicos. Iniciado a mediados de la década de 2000 por los Institutos Nacionales de Salud de EE. UU. (NIH), la Fase 1 del HMP fue diseñada para establecer una línea de base comprensiva de las comunidades microbianas en adultos sanos, enfocándose en cinco hábitats corporales principales: la cavidad oral, la piel, las narinas, el tracto gastrointestinal y el tracto urogenital. Esta fase generó extensos conjuntos de datos mediante la secuenciación del gen 16S rRNA, que proporcionó perfiles taxonómicos, y la metagenómica de escopeta del genoma completo, permitiendo obtener una visión funcional más profunda. Se tomó una muestra de un grupo de aproximadamente 300 individuos sanos en múltiples sitios del cuerpo, resultando en la recolección de más de 11,000 muestras físicas. De estas, casi 10,000 fueron analizadas por secuenciación de 16S rRNA, mientras que un subconjunto estratégico de aproximadamente 2,300 muestras se sometió a secuenciación metagenómica de escopeta. Este enfoque dual demostró que, si bien las composiciones de comunidades microbianas varían significativamente entre individuos, se identificaron vías metabólicas conservadas en comunidades dispares, un hallazgo que enfatizó la necesidad de estudiar la función microbiana además de la taxonomía. Los robustos conjuntos de datos de baseline (datos de individuos sanos usados como referencia posteriormente) generados durante la Fase 1 del HMP han proporcionado referencias invaluables para estudios comparativos posteriores que evalúan la disbiosis microbiana en varios estados de enfermedad \cite{aagaard2013human}.

Después de este trabajo fundamental, el Proyecto Integrativo del Microbioma Humano (iHMP) amplió el alcance al centrarse en contextos de enfermedades específicas. Durante un período de tres años, iHMP se enfocó en condiciones como el inicio de la diabetes tipo 2, la enfermedad inflamatoria intestinal y el parto prematuro relacionado con el embarazo \cite{integrative2014integrative}. A diferencia de su predecesor, iHMP empleó un diseño integrativo y longitudinal que combinaba datos multi-ómicos tanto del microbioma como del huésped. Este enfoque incorporó no solo análisis metagenómicos, sino también transcriptómica, proteómica y metabolómica junto con datos genómicos y epigenómicos del huésped, proporcionando así una visión dinámica y comprensiva de las interacciones huésped-microbio durante la progresión de la enfermedad. Los extensos conjuntos de datos producidos por ambas fases del HMP (superando los 48 terabytes y abarcando datos de más de 31,000 muestras de al menos 20 tipos diferentes de datos) han sido puestos a disposición de la comunidad científica a través del Centro de Coordinación de Datos del Proyecto del Microbioma Humano (HMPDACC) \cite{nih2019review}. El HMPDACC sirve como un repositorio centralizado que ofrece no solo acceso a datos de secuenciación en bruto, sino también a conjuntos de datos procesados que incluyen perfiles comunitarios, predicciones de genes, anotaciones funcionales y análisis de vías. El HMPDACC está construido sobre principios de datos FAIR (Findable, Accessible, Interoperable, and Reusable), asegurando que los investigadores de todo el mundo puedan utilizar los datos de manera efectiva. El portal ofrece capacidades avanzadas de búsqueda, herramientas interactivas de consulta y archivos de manifiesto descargables para facilitar la recuperación de datos tanto a granel como individuales \cite{creasy2021hmpdacc}.


\subsubsection{Bases de datos}

Gracias a la creciente aceptación del concepto FAIR en la comunidad científica, un número cada vez mayor de investigadores, tanto grandes como pequeños, están compartiendo los datos generados en sus proyectos. Este movimiento hacia la apertura facilita la accesibilidad, interoperabilidad y reutilización de los datos. Como respuesta, han surgido iniciativas internacionales que crean repositorios robustos para almacenar y gestionar esta información, tales como GenBank \cite{sayers2025genbank}, el European Nucleotide Archive \cite{o2025european} y el DNA Data Bank of Japan\cite{kodama2025ddbj}.

GenBank es un repositorio público integral de secuencias de nucleótidos y anotaciones biológicas administrado por el NCBI, parte del Instituto Nacional de Salud de Estados Unidos. GenBank es uno de los repositorios de secuencias de nucleótidos más completos, abarcando datos de diversas divisiones de organismos,incluyendo: secuencias de nucleótidos de ADN y ARN, genomas completos, y datos de transcriptómica y metagenómica. También incluye anotaciones funcionales y datos de polimorfismos, proporcionando información sobre variaciones genéticas y funciones de las secuencias.


El Archivo Europeo de Nucleótidos (ENA) es un repositorio público, gestionado por EMBL-EBI, que ofrece datos de secuenciación de todo tipo, permitiendo tanto la sumisión de datos como el acceso público a ellos. ENA recopila una amplia gama de datos relacionados con secuencias nucleotídicas, incluyendo secuencias genómicas completas, ADN y ARN, junto con datos de metagenómica y transcriptómica. Además, ofrece metadatos sobre los experimentos y proyectos de ensamblaje, así como lecturas de secuenciación cruda. 


El DDBJ (DNA Data Bank of Japan) es una base de datos que recoge, administrado por el National Institute of Genetics (NIG) en Japón y proporciona acceso a secuencias de ADN, ARN, secuencias genómicas completas y parciales. También ofrece datos de metagenómica, transcriptómica, anotaciones funcionales y lecturas de secuenciación cruda.


Actualmente estas tres organizaciones se alinean bajo la iniciativa de la Colaboración Internacional de Bases de Datos de Secuencias de Nucleótidos (INSDC). Este esfuerzo conjunto aborda varios objetivos significativos:

- Intercambio de Datos: Facilitan el intercambio continuo de información entre las organizaciones, asegurando que la información genómica sea accesible de manera uniforme a nivel global.

- Acceso Abierto: Ofrecen acceso libre y gratuito a datos de secuencias genéticas, fundamental para la investigación científica y el avance en tecnologías biológicas.

- Estandarización: Se centran en estandarizar formatos y protocolos para garantizar la interoperabilidad y coherencia internacional.

- Actualización Continua: Se comprometen a actualizar constantemente las bases de datos, integrando nuevas secuencias y descubrimientos científicos, reforzando así el progreso constante en el campo.




\textbf{WMS}

En el caso de las secuencias provenientes de WMS, a lo largo de la tesis, se han utilizado dos bases de datos (más concretamente, la misma pero dos versiones diferentes) WebOfLife (WoL) y WebOfLife2 (WoL2) \cite{zhu2019phylogenomics} para la asignación taxonómica de las muestras provenientes de WMS. Dichas bases de datos se con Bowtie2 \cite{langmead2012fast} para poder trabajar con ellas. De manera paralela se han preprocesado las muestras con fastp \cite{chen2018fastp} que nos permite realizar un gran numero de chequeos de calidad, siendo clave, el filtrado de lecturas de baja calidad y la eliminación de adaptadores. Una vez preprocesadas las muestras se alinean contra las bases de referencia indexadas usando Bowtie2. Una vez todas las muestras se hayan alineado, se utiliza Woltka \cite{zhu2022phylogeny} para la clasificación de los alineamientos y la asignación taxonómica en unidades genómicas operativas (OGUs), concepto homólogo a los ASVs de la secuenciación 16S.


%%%%%%%%%%%%%%%%%%%%%%%%%%%%%%%%%%%%%%%%%%%%%%%%%%%%%%%%%
%%%%%%%%%%%%%%%%%%%         Discussion              %%%%%%%%%%%%%%%%%%%%%%%
%%%%%%%%%%%%%%%%%%%%%%%%%%%%%%%%%%%%%%%%%%%%%%%%%%%%%%%%%


\chapter{Discusión}
\label{chap:discussion}

\chapterintro{
En conjunto los tres artículos que conforman esta Tesis Doctoral estudian, tanto el potencial del microbioma como fuente de biomarcadores, como la aplicación de técnicas de ML para el diagnóstico o prognóstico mediante el análisis de datos procedentes del microbioma. Siguiendo el hilo argumental de los tres artículos y los resultados obtenidos, esta sección aborda aspectos específicos y transversales generados a lo largo de la investigación.
}

La identificación de microorganismos o grupos de microorganismos que están relacionadas con la etiología o patogénesis de diversas enfermedades es una tarea altamente compleja. Esta requiere un enfoque multidisciplinar que integre técnicas avanzadas de microbiología, biología molecular y bioinformática para establecer conexiones claras y precisas. Por lo que, esta situación es el escenario perfecto para la búsqueda de patrones y relaciones que actúen como biomarcadores para el diagnóstico o prognóstico de dichas enfemedades. La motivación para el desarrollo de la presente tesis reside en afrontar y descubrir estas realciones usando técnicas computacionales.


Como ya hemos comentado, la caracterización microbiana por cultivo suponía un gran impedimento en el avance de este campo, principalmente por el reducido número de taxones que se podían cultivar y por el tiempo que tomaba realizar dichos cultivos. Pero, gracias al desarrollo de las nuevas tecnologías de secuenciación, y su abaratamiento a lo largo de los últimos años, se ha podido avanzar en la caracterización y el rol que desempeñan los innumerables microorganismos que conforman el microbioma humano.

De manera paralela, el avance de la bioinformática ha experimentado un crecimiento significativo en los últimos años. Este avance conlleva el desarrollo de herramientas capaces de procesar los datos crudos (resultado de la secuenciación) cada vez de manera mas rápida y eficaz. Así como, del aumento y anotación de bases de datos disponibles que nos permiten identificar, cada vez de manera más precisa, los taxones que se encuentran en el microbioma humano.

De la mano, en los ultimos años el número de investigaciones que utilizan estos perfiles microbiómicos para obtener posibles biomarcadores clínicos que puedan ayudar al diagnóstico, pronóstico y seguimiento de tratamiento en enfermedades ha aumentado. Un tema recurrente en muchos estudios es la aplicación de enfoques estadísticos sencillos, que incluyen medidas de la diversidad de la comunidad (diversidad alfa y beta), diferencias en la abundancia de taxones y la evaluación de la presencia o ausencia de especies.  


Estudios de alto impacto han utilizado con frecuencia métricas como el índice de Shannon y la Diversidad Filogenética de Faith para cuantificar la diversidad microbiana dentro de las muestras. Por ejemplo, estudios metagenómicos han demostrado que una disminución en la diversidad a menudo se correlaciona con la gravedad en una enfermedad, como se observa en pacientes con cáncer colorrectal o condiciones inflamatorias. Por el contrario, un alto nivel de diversidad alfa generalmente se asocia con un ecosistema intestinal estable y resiliente, lo que puede conferir protección contra el desarrollo de enfermedades. Los investigadores han demostrado empíricamente que incluso cambios sutiles en la diversidad alfa pueden servir como biomarcadores sólidos para fenotipos clínicos, especialmente cuando se combinan con medidas más refinadas de abundancia taxonómica \cite{yu2017metagenomic}.

La diversidad beta cuantifica las diferencias en la composición de comunidades microbianas entre muestras. Es especialmente útil cuando los investigadores buscan detectar diferencias sistemáticas entre cohortes de pacientes, como casos y controles o subtipos de enfermedades distintas. Los estudios han utilizado pruebas no paramétricas como PERMANOVA sobre métricas de distancia como la disimilitud de Bray-Curtis y las distancias UniFrac para evaluar si grupos de pacientes muestran diferencias estadísticamente significativas en sus estructuras generales de comunidades microbianas. Por ejemplo, estudios previos han informado que las diferencias en la diversidad beta pueden distinguir de manera confiable estados disbióticos de condiciones saludables, sentando las bases para el desarrollo de biomarcadores en oncología y enfermedades inflamatorias \cite{yilmaz2019microbial}. 

Como ya hemos comentado, además de las métricas de diversidad, se utilizan con frecuencia métodos estadísticos simples, como la comparación de abundancias relativas de taxones específicos, para evaluar la presencia o ausencia de especies microbianas que desempeñan roles críticos en la salud y la enfermedad. Estudios de alto impacto han demostrado que ciertos patógenos o especies clave presentan diferencias significativas de abundancia entre grupos de pacientes. Por ejemplo, varios estudios han identificado una sobre representación de \textit{Fusobacterium nucleatum} y un descenso de \textit{Faecalibacterium} en pacientes con cáncer colorrectal, hallazgos que han sido validados mediante comparaciones sencillas de abundancias de marcadores genéticos utilizando datos de secuenciación de 16S rRNA y metagenómica \cite{yu2017metagenomic, forbes2018comparative}.



Estos ejemplos demuestran la eficacia de las medidas de diversidad simples y el análisis de abundancias como biomarcadores clínicos iniciales, ofreciendo una visión sobre las alteraciones generales en la comunidad microbiana vinculadas a una enfermedad específica. Además, ambos enfoques facilitan la identificación de taxones diferenciales que más contribuyen a la separación entre grupos, destacando así potenciales biomarcadores candidatos. Por ello se decidió abordar la problemática del primer artículo \fcolorbox{Black}{Red}{[citar paper 1]}  de esta Tesis de manera similar. La estrategia seguida fue realizar estos estudios de diversidad microbiana a nivel de género, estudiar las abundancias relativas a nivel de filo y complementar el análisis con la comparación de abundancias de taxones entre los diferentes estados de condiciones patológicas. 

% Gracias a iniciativas (como HMP, AGP, DIABINMUNE, etc) y repositorios (como el ENA, GenBank o el DDBJ), se tiene acceso en todo el mundo a datos del microbioma relacionados con un número creciente de enfermedades. Además, con el paso de los años, se generan y almacenan datos de más pacientes, lo que ofrece la posibilidad de desarrollar metodologías capaces de obtener nueva información biológica en metacohortes de pacientes.


Los datos del microbioma destacan por su elevada dimensionalidad, su dispersidad, su composicionalidad y variabilidad interindividual sustancial. Estas caractersísticas íntrinsecas plantean importantes retos analíticos. Acabamos de ver como, aún así, los metodos estadísticos básicos son capaces de extraer información sustancial de los datos. Sin embargo estos métodos pueden tener problemas para captar interdependencias e interacciones complejas y no lineales entre taxones, la distribución desbalanceada de los taxones y la gran dispersión de los datos. En este contexto los metodos de ML han demostrado ventajas, manteniendo robustez frente a este tipo de estructura de datos, modelando relaciones complejas y mejorando las capacidades predictivas frente a otros enfoques.


Sin embrago, el uso de métodos de ML en este área (y en la biomedicina en general) presenta un gran obstáculo: el reducido tamaño muestral disponible en los estudios. Es decir, el número de muestras es mucho menor que el número de características (taxones y variables clínicas) de las que se disponen.

En el caso de estudios biomédicos de microbioma el tamaño muestral (generalizando) suele rondar entre los 50-100 pacientes, con mucha suerte. Mientras que si nos enfocamos en los taxones a estudiar, dependiendo del nival taxonómico del que hablemos, podemos movernos de los 10 a > 10.000 (de filo a especie/cepa). Y esto es un problema, ya que, a medida que exista un mayor número de características que de pacientes los modelos entrenados van a tender al sobreentrenamiento. Esto es, si pensamos en un problema de clasificación, a la hora de entrenar un modelo, este va a usar todas las caracteristicas para maximizar la diferencia entre las clases, sin embargo al agudizar tanto las diferencias encontradas, estas van a ser intrínsecas a los datos usados para entrenar los modelos. El resultado de esto es que el modelo va a tener una pobre capacidad de generalización, esto es, cuando se intenté utlizar para predecir en otras cohortes externas su rendimiento va a bajar drasticamente.

En este contexto, esta tesis, sigue un diseño experimental robusto, valiéndose de técnicas de selección de características para mantener solo aquellos taxones que suplan información a los modelos, y, eliminar aquellos redundantes o que puedan generar ruído en el análisis. Junto a estas técnicas el uso de una validacón cruzada anidada nos acerca a la construcción del modelo lo más generalizable posible.

Esta selección de características tiene un papel importante en el segundo artículo de la tesis \fcolorbox{Black}{Red}{[citar paper 2]}. Este artículo se vale de la selección de características para indetificar las especies más relevantes capaces de diferenciar las muestas entre sano y DT1. Los resultados muestran un aumento significativo en el rendimiento de los modelos si seleccionamos un subgrupo de especies del número total disponibles. La estrategia seguida fué la de utilizar una técnica de seleccion de caracteristicas de filtrado, lo que permite medir la relevancia de cada una de las especies con las clases de interés, identificando así cuantas y que especies utilizar para construir el modelo predictivo.


La potencia de las técnicas de ML no sólo se limitan a la construcción de modelos predictivos. Las técnicas de ML no supervisado se han convertido en una poderosa herramienta en biomedicina para descubrir agrupaciones naturales dentro de poblaciones heterogéneas de pacientes (sin necesidad de etiquetas preexistentes). Estas técnicas son especialmente adecuadas prar los análisis de datos ómicos de alta dimensionalidad (como es el caso del microbioma) ya que extraen automaticamente características latentes que pueden revelar subtipos o endotipos de enfermedades que no han sido descubiertos hasta la fecha. Esto es importante, ya que muchas enfermedades complejas muestran una variabilidad sustancial entre pacientes, y los enfoques estadísticos tradicionales o de ML supervisado que dependen de categorías clínicas predefinidas pueden pasar por alto dichas variaciones. Por lo que, las principales ventajas en estudios biomédicos de estos métodos no supervisados son su capacidad para interrogar los datos sin sesgo de etiquetas y que, superan las limitaciones de los pequeños tamaños muestrales en relación con el número de características medidas (lo que encaja perfectamente en biomedicina ya que las cohortes de pacientes suelen tener un tamaño limitado).

Estos métodos de ML no supervisado son la base para el desarrollo del tercer artículo \fcolorbox{Black}{Red}{[citar paper 3]}. Este artículo se vale de estás técnicas para la caracterización y extracción de subtipos de perfiles microbiómicos en el contexto de la VB. Los resultados muestran como estes subtipos se mantienen estables a lo largo de las distintas cohortes utilizadas en el artículo. La estrategia seguida fué la de utilizar varias técnicas de clusterización (con submuestreo y validación) hasta obtener un consenso del número de subgrupos óptimo a partir del perfil micorbiómico vaginal. 


En general, el uso de ML en el descubrimiento de biomarcadores del microbioma demuestra claras ventajas en términos de precisión predictiva y capacidad para gestionar datos heterogéneos y de alta dimensión, mientras que las técnicas estadísticas convencionales siguen proporcionando marcos confirmatorios esenciales para garantizar la interpretabilidad y una inferencia sólida. El éxito de estos enfoques integrados se pone de manifiesto en estudios en los que las canalizaciones basadas en ML han mejorado significativamente la precisión de la clasificación de enfermedades y la identificación de biomarcadores en comparación con los análisis basados únicamente en métodos convencionales.



Para concluir la sección de discusión, es crucial destacar las limitaciones identificadas a lo largo del desarrollo de esta tesis. Estas limitaciones se presentan a diferentes niveles: primero, en relación con cada uno de los artículos; segundo, de manera general de las diversas fuentes de variabilidad que pueden influir en los datos; y tercero, intrínseco al enfoque de composicional de estudiar el microbioma.

En el primer artículo una de las limitaciones es el reducido tamaño de la cohorte (79 pacientes), ya que, a mayor tamaño muestral mayor poder estadístico . Además se observa un desbalanceo significativo entre las clases de estudio de la diabetes (sano, prediabético y DT1), esto puede influir en la robustez y la habilidad de la metodología seguida para detectar diferencias significativas. Finalmente, se trata de una cohorte ubicada en una región geográfica determinada, podría no reflejar la heterogeneidad de poblaciones más amplias (limitando así la capacidad de generalización los resultados)

Respecto al segundo artículo, debido a la naturaleza del estudio (infantes con DT1) la validación de los modelos en cohortes externas no se ha podido llevar a cabo (inexistencia de cohortes del mismo tipo), por lo que no se puede evaluar la capacidad de generalización de los modelos construídos. Debido a que la variabilidad del microbioma entre infantes y adultos es tan alta no se ha podido validar en cohortes extrernas de pacietnes adultos con T1D.

En lo relativo al tercer artículo, este presenta la limitación de que, aunque se han utilizado varias cohortes, la distribución geográfica de las mismas está bastante acotada, por lo que, las puede limitar la capacidad de generalización de los resultados obtenidos. Además, en la cohorte discovery la asignación taxonómica se ha realizado basándose en OTUs (que como ya hemos comentado, no es la estrategia más recomendada actualmente), mientras que en el resto de cohortes se ha realizado basándose en ASVs, por lo que puede existir algún sesgo, Adicionalemente, aunque k-means y glmnet son usados extensivamente y muestran robusatez, pueden presentar pequeños sesgos internos que pueden influir en el perfilado microbiómico.

De manera más general,una grna limitación son las fuentes de variablilidad con la que los estudios de este tipo tienen que lidiar y que pueden acabar afectando a los downstream análisis. Ya en diseño experimental, donde se eligen que regiones del 16S hay que secuenciar existe variablidad, la cual, se ve acentudada debido al amplio catalógo de kits de extracción del ADN. A mayores existe mucha variablidad derivada de la eleccion de utilizar uan determinada plataforma de secuenciación, así como de la estreategia llevada a cabo (Ilumina, PacBIO, Ion Torrent, etc). Finalmente, otro de los puntos clave donde se introduce gran variabildad es la asignación taxonómica de los datos, donde influyen las pipelines utilizadas (DADA2, Qiime2, MOTHUR, etc), las bases de datos utilizadas (Greengenes, SILVA, RDP, etc) y la estrategia seguida (OTUs o ASVs). Aunque el ML es bueno manejando esta heterogeneidad, inevitablemente se añaden sesgos a los análisis.

Finalmente, la ultima limitación a destacar es inherente al estudio composicional del microbioma, esto es, al final los estudios que se están realizando se basan en datos procedentes de una instantátena en un momento dado del perfil microbiómico de los pacientes. Intentar reducir la complejidad biológica a un fotograma determinado es simplificar excesivamente procesos dinámicos y multifacéticos de cualquier enfermedad.
Por lo que añadir más puntos temporales (estudios longitudinales) o datos procedentes de otras fuentes ómicas (metatranscriptoma, metaproteoma metaboloma, etc) ayudaría a capturar la complejidad e información biológica de manera más realista/representativa.



% Diverisda ALFA

La medida complementaria (1 – D) o su recíproco se utilizan a menudo para proporcionar una interpretación más intuitiva, donde valores más altos indican una mayor diversidad. El índice de Simpson otorga mayor énfasis a las especies abundantes, dado que al elevar al cuadrado las proporciones se amplifica la influencia de los taxones comunes en relación con los raros, lo que lo hace menos sensible a la presencia de especies raras en comparación con el índice de Shannon. Esta característica del índice de Simpson puede ser particularmente ventajosa cuando el objetivo es evaluar la estructura de dominancia dentro de una comunidad microbiana, ya que puede revelar escenarios en los cuales unas pocas especies predominan o cuando varios taxones contribuyen de forma más equitativa al perfil general. Sin embargo, aunque el índice de Simpson es robusto frente a errores de muestreo y es más fácil de interpretar cuando se centra en la dominancia de especies, puede subrepresentar la contribución de especies raras o de baja abundancia que podrían ser ecológicamente significativas \cite{finotello2018measuring}.

En conjunto, los índices de Shannon y Simpson encapsulan visiones complementarias de la diversidad microbiana. Mientras que el índice de Shannon proporciona un resumen integral que captura tanto la riqueza como la uniformidad, el índice de Simpson ofrece una perspectiva que subraya la uniformidad de las especies altamente abundantes, al tiempo que minimiza el peso de los taxones raros. Su aplicación conjunta permite a los investigadores desmenuzar las complejas capas de la estructura de la comunidad microbiana y comprender las implicaciones tanto de la diversidad general como de las jerarquías de dominancia dentro de las comunidades \cite{kleine2021data}.


% Diversidad BETA


Técnicas de ordenación como el análisis de coordenadas principales (PCoA) o el escalado multidimensional no métrico (NMDS) son típicamente empleadas para visualizar los resultados de la diversidad \textit{beta} en un espacio de dimensionalidad reducida, facilitando así la detección de patrones de agrupamiento que corresponden a datos o variables clínicas \cite{knight2018best}. Además, herramientas estadísticas como el Análisis de Varianza Multivariado por Permutación (PERMANOVA) y el Análisis de Similitudes (ANOSIM) se utilizan para probar rigurosamente la significancia de las diferencias en la diversidad \textit{beta} entre grupos predefinidos. Estos métodos multivariados proporcionan valores p que indican si las diferencias observadas en la composición comunitaria son mayores de lo que se podría esperar por muestreo aleatorio, reforzando así las conclusiones sobre las firmas microbianas específicas de cada grupo. Es importante notar que los análisis de diversidad \textit{beta} requieren una cuidadosa normalización de los datos de secuenciación para controlar las diferencias en la profundidad de secuenciación, así como prestar atención a la composicionalidad inherente de los datos del microbioma, que puede influir en las mediciones de distancia \cite{knight2018best}.


% Machine Learning no supervisado
% Distancia de Mahalanobis
\textbf{Distancia de Mahalanobis:} mide la distancia entre un punto y una distribución, teniendo en cuenta la covarianza entre las variables.
\begin{equation}
d_{\text{mahalanobis}}(\mathbf{x}, \mathbf{y}) = \sqrt{(\mathbf{x} - \mathbf{y})^T \mathbf{\sum}^{-1} (\mathbf{x} - \mathbf{y})}
\end{equation}

\begin{itemize}
\item Donde:

\item $x$ y $y$ son vectores que representa un punto en el espacio (n)-dimensional.
\item $\sum$ es la matriz de covarianza de los datos.
\end{itemize} 

% Enlaces
\begin{itemize}
\item \textbf{Enlace simple:} también conocido como distancia mínima, calcula la distancia entre dos \textit{clusters} como la mínima distancia entre cualquier par de puntos que pertenecen a \textit{clusters} diferentes

\item \textbf{Enlace completo:} o distancia máxima, calcula la distancia entre dos \textit{clusters} como la máxima distancia entre cualquier par de puntos, uno de cada \textit{cluster}.

\item \textbf{Enlace promedio:} calcula la distancia entre dos \textit{clusters} como el promedio de todas las distancias entre los puntos de ambos \textit{clusters}.

\item \textbf{Método de Ward:} minimiza la suma total de las varianzas dentro de cada \textit{cluster} al fusionar dos \textit{clusters}. Este método intenta mantener la homogeneidad dentro de los \textit{clusters}, haciéndolos lo más compactos posible.
\end{itemize} 


\paragraph{MOFA}

MOFA es una técnica integrativa basada en modelos bayesianos de reducción de dimensionalidad diseñada para el análisis simultáneo de diferentes tipos de datos ómicos (transcriptoma, proteoma, micorbioma, etc.). El funcionamiento de MOFA (representado en la Figura \ref{fig:MOFA}) se basa en la descomposición de cada matriz de datos ómico $Y_{m}$ en dos componentes principales: una matriz de factores latentes $Z$, que captura la actividad subyacente de cada muestra, y una matriz de pesos o cargas $W_{m}$ asociada a cada tipo de dato ómico, la cual indica la importancia de cada característica. La formulación matemática del modelo establece que cada entrada $y$ en la matriz $Y_{m}$ se puede aproximar por la suma ponderada de los factores latentes. Para garantizar interpretabilidad y evitar el sobreajuste, MOFA impone restricciones de dispersión en las matrices de pesos, de modo que solo un subconjunto pequeño de características tenga una contribución significativa en cada factor. Esta estrategia favorece la identificación de patrones moleculares específicos y facilita la asignación de significado biológico a los factores latentes.

Cuando MOFA se aplica en el contexto de un solo conjunto de datos, como ocurre en estudios de microbioma, el método esencialmente se reduce a un proceso de análisis de reducción de dimensionalidad comparable a un PCA, aunque enriquecido por la estructura bayesiana y las restricciones de dispersión que implementa. En este escenario, la matriz de datos que representa las abundancias microbianas se factoriza en una matriz de factores latentes y una matriz de pesos, donde los factores latentes resumen las principales fuentes de variación en la comunidad microbiana. Aunque en un análisis multiómico MOFA puede distinguir entre factores compartidos y específicos de cada tipo ómico, en un solo tipo de dato la estructura resultante refleja las variaciones dominantes dentro de la misma población microbiana, posibilitando la identificación de variables clínicamente relevantes. Además, la robustez del modelo frente la naturaleza dispersa e inflación de ceros de los datos de microbioma, confiere una ventaja notable frente a métodos tradicionales. En este contexto, MOFA actúa como una extensión poderosa de los métodos clásicos de reducción de dimensionalidad, proporcionando además mecanismos de regularización y una interpretación probabilística de los factores que pueden mejorar la identificación de firmas microbianas asociadas a estados de salud o enfermedad.

\begin{figure}[h]
    \centering
    \includegraphics[width=\textwidth]{img/introduction/17_MOFA.png}
    \caption{\textbf{Diagrama de MOFA.} Cada matriz representa un tipo de datos. MOFA encuentra factores comunes entre estas matrices, revelando patrones compartidos y diferencias específicas.  Imagen creada con BioRender \cite{biorender2025}.}
    \label{fig:MOFA}
\end{figure}



%%%%%%%%%%%%%%%%%%%%%%%%%%%%%%%%%%%%%%%%%%%%%%%          DISCUSION                %%%%%%%%%%%%%%%%%%%%%%%%%%%%%%%%%%%%%%%%%%%%%%%%%%%%%%%%

\chapter{Discusión}
\label{chap:discussion}

\chapterintro{
En conjunto los tres artículos que conforman esta Tesis Doctoral estudian, tanto el potencial del microbioma como fuente de biomarcadores, como la aplicación de técnicas de ML para el análisis de datos del microbioma con el fin de apoyar procesos diagnósticos. Siguiendo el hilo argumental de los tres artículos y los resultados obtenidos, esta sección aborda aspectos específicos y transversales generados a lo largo de la investigación.
}

\section{Introducción general de la discusión}

\textbf{Breve recapitulación del objetivo principal de la tesis}

El objetivo principal de esta tesis fue explorar el potencial del microbioma humano como fuente de futuros biomarcadores clínicos, utilizando técnicas computacionales. A través de un análisis detallado de datos metagenómicos, buscamos identificar patrones específicos que puedan correlacionarse con diversas condiciones de salud. Asimismo, crear modelos predicitivos a partir de estos patrones específicos que nos permita clasificar pacientes de una determinada patología. Este enfoque no solo pretende ampliar nuestro conocimiento sobre la relación entre el microbioma y el estado de salud, sino también establecer posibles vías alternativas para futuras investigaciones en medicina.

\textbf{Relevancia del tema}

El estudio del microbioma humano ha cobrado gran relevancia en los últimos años debido a su potencial como fuente de biomarcadores para enfermedades complejas. Como se ha expuesto en la presente Tesis Doctoral, cada vez hay más evidencias que defienden que el microbioma modula funciones esenciales del sistema inmune, el metabolismo y la respuesta inflamatoria. Esto se refleja en investigaciones que han asociado alteraciones en el microbioma con patologías como la DT2, algunos tipos de cáncer y trastornos autoinmunitarios \cite{schirmer2016linking, fan2021gut, bakir2022inflammatory}. Por lo que, los biomarcadores microbianos pueden ofrecer una ventana para la detección temprana y un entendimiento más profundo de estas enfermedades. Además, el estudio en profundidad de estos biomarcadores y su comparación con los ya existentes en el estado del arte puede llevar a mejoras en las estrategias diagnósticas o terapéuticas. 


\textbf{Justificación metodológica general}

El enfoque metodológico adoptado en la presente Tesis Doctoral se fundamenta en la integración de técnicas estadísticas de análisis de abundancia diferencial, técnicas de selección de características y técnicas de ML para abordar las particularidades y limitaciones inherentes al análisis del microbioma humano. Por una parte, los métodos de abundancia diferencial específicos para microbiomas (al utilizar transformaciones logarítmicas y modelos estadísticos adaptados a la sobredispersión e inflación de ceros) aseguran que las inferencias realizadas sean consistentes con las propiedades intrínsecas de los datos. Por otra parte, la utilización de estrategias de selección de características posibilita la reducción drástica de la dimensionalidad del conjunto de datos sin sacrificar la información crítica, facilitando la extracción de biomarcadores que exhiben alta capacidad discriminativa y relevancia biológica. Finalmente, la estrategia de ML seguida permite construir modelos predictivos capaces de manejar la elevada variabilidad en las comunidades microbianas entre individuos posibilitando así la identificación de firmas microbianas que pueden predecir estados patológicos con elevado poder de clasificación y generalización.


\section{Evaluación crítica transversal de los resultados}

\textbf{Enfoque Data-driven frente a Knowledge-driven}

La motivación para utilizar técnicas basadas en datos (\textit{Data-driven}) en esta Tesis Doctoral surge de la complejidad intrínseca, alta dimensionalidad y dispersión de los datos del microbioma. Estas características generan desafíos significativos al intentar entender las interacciones entre hospedadores y microorganismos, así como al identificar posibles biomarcadores para uso clínico. Los enfoques basados en el conocimiento (\textit{Knowledge-driven}) a menudo dependen de hipótesis biológicas, bases de datos curadas y modelos mecanicistas. Sin embargo, están limitados por el conocimiento biológico actual y frecuentemente no logran captar las complejas, no lineales y multidimensionales asociaciones presentes en los datos del microbioma. Por otro lado, los métodos dirigidos por datos operan sin depender de hipótesis o vías biológicas preexistentes, lo que es ventajoso en sistemas multifactoriales como el microbioma intestinal \cite{dakal2025advanced}. Además, los enfoques basados en ML permiten analizar directamente los datos en bruto, descubriendo patrones ocultos que quedarían invisibles en un marco basado en el conocimiento \cite{dreisbach2023comparison}. Esto los hace especialmente atractivos para el descubrimiento exploratorio de biomarcadores, ya que pueden revelar asociaciones y patrones no detectables por métodos estadísticos tradicionales \cite{hemmati2024using}.


En el estado del arte actual se utilizan muchos enfoques basados en el conocimiento en el campo del microbioma \cite{chong2017computational, bauer2018network,  luo2024progress}. Sin embargo no es un enfoque recomendable a la hora de buscar nuevos potenciales biomarcadores al estar restringido a hipótesis biológicas preexistentes. No obstante es una buena estrategia complementaria para confirmar o reafirmar posibles biomarcadores (al introducirlos en una red de conocimiento ya existente) encontrados con enfoques basados en datos. Este enfoque nos ha permitido identificar taxones que inicialmente no se relacionaban ni se conocía su papel en diversas patologías. A lo largo de los tres trabajos \cite{paper1, paper2, paper3}, al analizar datos brutos en lugar de un subconjunto basado en conocimiento previo, hemos descubierto taxones ya mencionados en la literatura, además de otros que eran desconocidos o poco documentados. Esto abre la puerta a investigar el papel de estos taxones en el desarrollo de enfermedades y su potencial como biomarcadores.

\textbf{Composicionalidad}

En el análisis del microbioma, tras estudiar las ventajas y contras de los distintos tipos de normalizaciones y transformaciones \cite{gloor2016compositional, gloor2017microbiome, hong2022rarefy}, hemos optado por utilizar transformaciones composicionales para abordar la problemática de la interdependencia entre taxones en los datos procedentes del microbioma. Esta característica intrínseca, que implica que la presencia de un taxón está directamente relacionada con la presencia de otros, puede conducir a interpretaciones erróneas si se analizan los datos en crudo. Por tal motivo, se han implementado transformaciones como el TSS y la CLR. Estas transformaciones buscan normalizar los datos y minimizar las distorsiones provocadas por la composicionalidad, permitiendo una representación más precisa de las relaciones entre los diferentes taxones y mejorar las comparaciones entre estudios. 

Actualmente en la literatura todavía se sigue usando metodos de normalización para datos porcedentes de 16S rRNA que no tienen en cuenta la composicionalidad de los datos. Numerosos trabajos siguen abogando por la rarefacción a un numero fijo de lecturas, tenemos ejemplo tanto en diabetes \cite{lambeth2015composition, dash2023functional, siptroth2023variation} como en BV \cite{chen2018vaginal, onywera2021predictive} lo que va a introducir un sesgo, en mayor o menor medida, a los resultados obtenidos. Por ello, tras el desarrollo de esta Tesis Doctoral, recomendamos el uso de las transformaciones logarítmicas previamente mencionadas para manejar adecuadamente las características del microbioma y favorecer la reproducibilidad de los resultados. Gracias a dichas transformaciones en el primer trabajo \cite{paper1} se ha utilizado un método como LinDa y se extrajeron patrones significativos. En el caso de los otros trabajos \cite{paper2, paper3}, se comparó el rendimiento aplicando estas transformaciones frente a los datos crudos, y se observó que el rendimiento se reducía al no aplicar dichas transformaciones (se puede ver en la figura 3 de los datos suplementarios del trabajo \cite{paper3}). Además, en este último trabajo aplicar estas transformaciones ayuda a reducir la variabilidad entre los diferentes enfoques utilizados (OTUs vs ASVs) a lo largo de las distintas cohortes.

\textbf{Dimensionalidad}

Para afrontar el desafío de la alta dimensionalidad en los datos del microbioma, se ha optado por realizar filtrados rigurosos (conteos totales y prevalencia de taxones) a la hora de extraer la tabla de taxones y técnicas de selección de características que permiten reducir el número de taxones considerados en la construcción de modelos de ML. Las cuales han demostrado ser efectivas en este contexto \cite{segata2011metagenomic}. Este enfoque no solo optimiza los modelos al disminuir el ruído y las variables redundantes, sino que también mejora su interpretabilidad y rendimiento. 

Actualmente el problema sobre la dimensionalidad está bastante arraigado en el estudio de datos ómicos, incluyendo el microbioma. La mayoría de artículos del estado del arte que utilizan enfoques basados en ML suelen realizar algún tipo de selección de características. Ya sea a través de algoritmos que realizan su propia selección (como pueden ser RF o glmnet) o por algún tipo de selección de características que clasifique dichas características según su importancia \cite{beck2015machine, biassoni2020gut, bakir2021discovering, bakir2022inflammatory, chavarria2025using}. Esta selección de características tiene un papel importante en el segundo trabajo de esta Tesis Doctoral \cite{paper2}. Este artículo se vale de la selección de características para identificar las especies más relevantes capaces de diferenciar las muestras entre sano y DT1. Los resultados muestran un aumento significativo en el rendimiento de los modelos si seleccionamos un subgrupo de especies del número total disponibles. La estrategia seguida fué utilizar una técnica de selección de características de filtrado, lo que permite medir la relevancia de cada una de las especies con las clases de interés, identificando así cuantas y que especies utilizar para construir el modelo predictivo.


\textbf{Variabilidad}

La variabilidad inter-muestras constituye otro obstáculo significativo en el análisis del microbioma, un reto que hemos abordado mediante el uso de técnicas avanzadas de validación. En el caso del aprendizaje no supervisado el uso del enfoque CCP nos ha permitido la construcción de subgrupos robustos y generalizables. Gracias al submuestro, la evaluación de los algoritmos de \textit{clustering} y la evaluación de los \textit{clusters} producidos de manera consensuada, nos aseguramos que el modelo final de \textit{clustering} obtenido es robusto y no presenta sobreentrenanmiento a un subgrupo de los datos originales. En el caso del aprendizaje supervisado al realizar una NCV nos garantiza un modelo robusto y generalizable que contemple la variabilidad inherente a las muestras biológicas \cite{varma2006bias}. Al segmentar el conjunto de datos en múltiples subconjuntos de entrenamiento y prueba, se logra una estimación fiable del rendimiento del modelo y se minimiza el riesgo de sobreajuste. 

Esto se ve reflejado en los trabajos \cite{paper2, paper3}, donde gracias a esta validación se ha podido manejar la diversidad en los datos de entrada, habilitando la construcción de modelos con poca dispersión en su rendimiento (independientemente del pool de entradas), lo que se traduce en un modelo con altas capacidades de clasificar pacientes entre sanos y DT1 en el trabajo \cite{paper2} y un modelo en el trabajo \cite{paper3} que es capaz de predecir los VCSs incluso cuando los datos de entrada provienen de fuentes heterogéneas.

\textbf{Dispersión}

Para mitigar la alta dispersión presente en los datos del microbioma, hemos promovido el uso de algoritmos de ML que han demostrado ser altamente eficaces en este contexto. Los RF, al generar y combinar múltiples árboles de decisión, permiten reducir la varianza y mejorar la estabilidad del modelo, siendo efectivos frente a la dispersión. Asimismo, los Gradient Boosting Machines, como XGBoost, se han revelado robustos ante outliers y datos ruidosos, gracias a su optimización iterativa con regularización, facilitando así el manejo de datos dispersos. Estos algoritmos no solo gestionan la dispersión intrínseca de los datos microbianos, sino que también mejoran la precisión y generalización de los modelos predictivos, contribuyendo así al desarrollo de biomarcadores consistentes en el análisis del microbioma.

El uso de estos algoritmos ha sido clave en los trabajos \cite{paper2, paper3}. Gracias al uso de algoritmos que se regularizan internamente se ha evitado que taxones que no tienen importancia o puedan confundir a los modelos mediante la penalización de los mismos. 

\textbf{OTUs vs ASVs?}

Uno de los principios centrales de la investigación sobre el microbioma es la elección de la unidad bioinformática para la inferencia de secuencias, siendo las OTUs dependientes de umbrales de agrupamiento fijos, mientras que las ASVs han emergido como el método preferido para lograr mayor resolución y reproducibilidad. En la introducción, se comentó que los métodos basados en OTUs normalmente agrupan secuencias en un umbral de similitud del 97\%, lo que puede enmascarar variaciones sutiles pero biológicamente significativas, mientras que los enfoques basados en ASVs usan los datos de secuencia en bruto para inferir variantes de secuencia exactas que difieren incluso en un nucleótido. En esta investigación hemos explorado la utilidad de las OTUs y las ASVs para el análisis del microbioma. Nuestros resultados indican que ambos enfoques son útiles según los contextos evaluados en los trabajos \cite{paper1, paper2, paper3}.

En el primer trabajo \cite{paper1}, se aplicó un flujo de trabajo basado en ASVs para realizar un análisis de abundancia diferencial entre sujetos sanos y pacientes con DT2, lo que resultó en la identificación de taxones específicos con diferencias significativas de abundancia entre los grupos. La ventaja de utilizar un enfoque basado en ASV en este análisis radica en que preserva la variación secuencial a escala fina que podría perderse con la agrupación de OTUs, permitiendo así detectar diferencias ecológicas sutiles entre los grupos de pacientes.

En contraste, el segundo trabajo \cite{2} se basó en un análisis de OTUs para extraer una firma metagenómica que se ingresó en un clasificador RF con el fin de distinguir entre pacientes sanos y enfermos de DT1. Aunque la agrupación de OTUs tiende a fusionar secuencias estrechamente relacionadas en \textit{clusters} más amplios y puede inflar las estimaciones de riqueza, numerosos estudios han demostrado que los métodos basados en OTUs aún capturan patrones ecológicos robustos a niveles taxonómicos superiores \cite{liu2023independent}. De hecho, en este segundo trabajo, el uso de OTUs permitió el desarrollo de un modelo de clasificación que generó una firma metagenómica altamente predictiva, sugiriendo que la resolución taxonómica proporcionada por la agrupación de OTUs era suficiente para los fines de clasificación clínica.

El tercer trabajo \cite{paper3} se centró en el \textit{clustering} del microbioma vaginal utilizando OTUs y se validó el modelo de \textit{clustering} en cohortes independientes con procesamiento basado tanto en OTUs como en ASVs. En este trabajo, el modelo de clustering derivado de datos de OTUs resultó ser lo suficientemente robusto para ser validado en muestras procesadas mediante un enfoque de ASVs, reflejando la observación de que los principales patrones biológicos a nivel comunitario se conservan en gran medida entre ambos enfoques. Este hecho respalda hallazgos previos que indican que, aunque las estimaciones de riqueza absoluta y las asignaciones taxonómicas a escala fina pueden diferir significativamente, los patrones generales de beta diversidad, la topología de \textit{clustering} y las señales biológicas son a menudo similares al comparar conjuntos de datos basados en OTUs y ASVs \cite{chiarello2022ranking}. La convergencia de los resultados de \textit{clustering} a través de múltiples cohortes sugiere que los \textit{clusters} basados en OTUs capturan la estructura comunitaria básica, aunque pueden subrepresentar la microdiversidad que los análisis basados en ASVs proporcionan.

En conclusión, las discusiones presentadas en estos tres artículos destacan colectivamente que, si bien los enfoques basados en OTUs y ASVs difieren en sus fundamentos teóricos y en sus resultados prácticos, ambas pueden generar resultados robustos y biológicamente significativos cuando se aplican con una atención cuidadosa a los protocolos de preprocesamiento y filtrado. La mayor resolución de los ASVs resulta especialmente valiosa cuando el objetivo de la investigación es detectar diferencias sutiles en la abundancia entre grupos de estudio (y más si se buscan especies singulares o incluso cepas), como se evidencia en \cite{paper1}. Por otro lado, la robustez y la relativa simplicidad del agrupamiento en OTUs han demostrado ser suficientes para extraer firmas metagenómicas predictivas para la clasificación de enfermedades, como se observa en \cite{paper2}, y para capturar la agrupación a nivel comunitario en \cite{paper3}. En última instancia, estos hallazgos refuerzan la idea de que la elección del método debe alinearse con los objetivos del estudio, reconociendo al mismo tiempo que el rendimiento analítico posterior (ya sea en el análisis de abundancia diferencial, la clasificación mediante modelos o \textit{clustering}) permanecerá en gran medida consistente a nivel macroecológico una vez que se apliquen rigurosas medidas de control de calidad.


\textbf{Modelos lineales vs modelos no lineales?} 

En el segundo trabajo \cite{paper2}, la tarea de predecir pacientes DT1 utilizando datos del microbioma intestinal presenta un desafío único, ya que el ecosistema intestinal es altamente complejo, con interacciones microbiales que son tanto no lineales como multifactoriales en naturaleza. En este contexto, el microbioma intestinal refleja una amplia gama de taxones en interacción cuyas relaciones son moduladas por la colonización del nicho, factores ambientales abióticos y respuestas inmunitarias del huésped. Los algoritmos de ML no lineales, como RF, están inherentemente equipados para capturar estas complejas estructuras de dependencia, ya que operan mediante la construcción de un conjunto de árboles de decisión que dividen jerárquicamente el espacio de características. Este enfoque en conjunto permite a los modelos RF identificar interacciones y asociaciones de orden superior entre características microbianas que pueden ser críticas para discriminar entre estados de enfermedad y condiciones saludables.

Por otro lado, el tercer trabajo \cite{paper3} se centra en la estratificación de pacientes utilizando perfiles del microbioma vaginal y emplea un modelo lineal, como es glmnet. Como se ha comentado el microbioma vaginal se caracteriza por ser mucho menos heterogéneo, presentar una reducidad diversidad y donde las comunidades microbianas tienden a estar dominadas por especies de \textit{Lactobacillus}. En tales casos, los límites que separan los perfiles comunitarios asociados con diferentes grupos de pacientes parecen ser más linealmente separables, lo que permite que modelos lineales como glmnet capturen efectivamente estas diferencias con el beneficio adicional de la interpretabilidad del modelo. La claridad en la delimitación de la comunidad microbiana en el nicho vaginal con frecuencia conduce a una situación en la que la suposición más simple de efectos aditivos es suficiente para la estratificación de pacientes, reduciendo así la necesidad de modelado no lineal más complejo.

En conclusión, el rendimiento superior de algoritmos no lineales como el RF en el contexto de la predicción de DT1 basada en el microbioma intestinal \cite{paper2} se debe principalmente a la necesidad de modelar interacciones intrincadas, de alta dimensión y no lineales inherentes al ecosistema microbiano intestinal, las cuales se ven agravadas por una alta variabilidad interindividual y complejas interacciones huésped-microbio. En contraste, el uso efectivo de un modelo lineal, como glmnet, en la estratificación de pacientes en función de los perfiles del microbioma vaginal \cite{paper3} refleja la naturaleza relativamente estructurada y linealmente separable de las comunidades microbianas vaginales, donde se pueden delimitar características clave con precisión, aprovechando los beneficios de la interpretabilidad y la escasez del modelo. Las diferencias en el contexto de la enfermedad y la heterogeneidad de los nichos microbianos son, por lo tanto, fundamentales para explicar por qué los métodos no lineales superan en un escenario, mientras que los modelos lineales son más efectivos en otro. Para concluir, esta evidencia sugiere que el contexto de la enfermedad, la heterogeneidad del nicho microbiano, la dimensionalidad de las características y la estructura general de los datos influyen en el rendimiento relativo de los modelos de ML lineales frente a los no lineales en la investigación del microbioma.


\textbf{Microbioma longitudinal}

El microbioma humano es un ecosistema altamente dinámico que exhibe una variabilidad temporal considerable tanto a nivel composicional como funcional, una característica que ha sido reconocida en la literatura y que plantea desafíos para la clasificación confiable de enfermedades utilizando datos del microbioma \cite{aldars2021systematic}. Nuestros trabajos \cite{paper2, paper3}, al igual que el de otros \cite{eck2017robust}, han demostrado que incluso cuando se utiliza solo una muestra por paciente para la clasificación de enfermedades, se pueden lograr resultados robustos, lo que sugiere que la señal microbiana relacionada con la patología puede ser lo suficientemente resistente como para emerger a pesar de las fluctuaciones temporales subyacentes. Es importante destacar que, aunque el microbioma experimenta cambios pronunciados a lo largo de horas, días y semanas, muchos estudios han revelado que las características centrales que proporcionan utilidad diagnóstica son, de hecho, más estables de lo que se podría esperar \cite{mehta2018stability}.

A la vista de la robustez de los resultados obtenidos, a pesar del comportamiento variable del microbioma, indican que los modelos que empleamos probablemente están aprovechando características que representan ya sea una respuesta promedio o aspectos altamente conservados del microbiota que son menos sensibles a las perturbaciones temporales. En esencia, aunque una única instantánea puede no captar todo el espectro de fluctuaciones temporales, parece encapsular un estado basal que es suficientemente representativo de las firmas microbianas relevantes para la enfermedad. Esta observación concuerda con la evidencia reportada en varios estudios, que sugieren que la robustez de la clasificación de enfermedades no depende únicamente de capturar cada fluctuación transitoria, sino más bien de detectar los marcadores microbianos persistentes y dominantes que diferencian la salud de la enfermedad.

%\textbf{Profundidad taxonómica}




\section{Evaluación crítica de los artículos}

- Discutir los principales resultados de cada uno de tus artículos. Piensa también en escribir un párrafo por paper. Importante también citar bibliografía y comparar tus resultados con otros estudios


- En el primer trabajo \cite{paper1}, el enfoque seguido nos ha permitido identificar taxones que podrían jugar un papel protector en SM o DT2 que no han sido estudiados ni reportados previamente en este contexto. En el caso de SM, se ha relacionado la especie \textit{Parabacteroides merdae} con los controles sanos. Dicha especie no se ha estudiado en mucha profundidad, pero un estudio muy reciente realizado en ratones ha expuesto que puede tener un rol beneficioso a través del catabolismo de aminoácidos de cadena ramificada (este tipo de aminoácido se relaciona con el aumento del riesgo de enfermedades metabólicas) \cite{qiao2022gut}. Por otro lado, en el caso de DT2, cuando se comparan los pacientes prediabéticos o diabéticos frente a los sanos, se observa una disminución en la proporción de dos géneros de los que no se tiene mucho conocimiento: \textit{Lachnospiraceae GCA-900066575} y \textit{Defluviitaleaceae UCG-011}. Al igual que otros miembros del orden \textit{Clostridiales}, podrían participar en la fermentación de fibras y la producción de AGCCs, que, como ya vimos, tienen efectos antiinflamatorios y reguladores sobre la función de la barrera intestinal. Además, estudios muy recientes realizados en ratones señalan que \textit{GCA-900066575} se relaciona negativamente con niveles de glucosa en sangre en ayunas y la resistencia a la insulina \cite{liu2023vitamin}, así como mejoras en la enfermedad del hígado graso no alcohólica \cite{qu2023akkermansia, guo2023protective}


- \textcolor{red}{En el segundo trabajo \cite{paper2}}



- \textcolor{red}{En el tercer trabajo \cite{paper3}}



\section{Limitaciones del estudio}

Para concluir la sección de discusión, es crucial destacar las limitaciones identificadas a lo largo del desarrollo de esta tesis. Estas limitaciones se presentan a diferentes niveles: primero, en relación con cada uno de los artículos; segundo, de manera general de las diversas fuentes de variabilidad que pueden influir en los datos; y tercero, intrínseco al estudio del microbioma.

En el primer trabajo \cite{paper1} una de las limitaciones es el reducido tamaño de la cohorte (79 pacientes), ya que, a mayor tamaño muestral mayor poder estadístico. Además, se observa un desbalanceo significativo entre las clases de estudio de la diabetes (sano, prediabético y DT1), esto puede influir en la robustez y la habilidad de la metodología seguida para detectar diferencias significativas. Finalmente, se trata de una cohorte ubicada en una región geográfica determinada, podría no reflejar la heterogeneidad de poblaciones más amplias (limitando así la capacidad de generalización los resultados).

Respecto al segundo trabajo \cite{paper2}, debido a la naturaleza del estudio (infantes con DT1) la validación de los modelos en cohortes externas no se ha podido llevar a cabo (inexistencia de cohortes del mismo tipo), por lo que no se puede evaluar la capacidad de generalización de los modelos construídos. Debido a que la variabilidad del microbioma entre infantes y adultos es tan alta no se ha podido validar en cohortes externas de pacientes adultos con T1D.

En lo relativo al tercer trabajo \cite{paper3}, este presenta la limitación de que, aunque se han utilizado varias cohortes, la distribución geográfica de las mismas está bastante acotada, por lo que las puede limitar la capacidad de generalización de los resultados obtenidos. Además, en la cohorte basal la asignación taxonómica se ha realizado basándose en OTUs (que como ya se ha comentado, no es la estrategia más recomendada actualmente), mientras que en el resto de cohortes se ha realizado basándose en ASVs, por lo que puede existir algún sesgo metodológico, Adicionalmente, aunque k-means y glmnet son usados extensivamente y muestran robustez, pueden presentar pequeños sesgos internos que pueden influir en el perfilado microbiómico.

De manera más general, una gran limitación son las fuentes de variabililidad con la que los estudios de este tipo tienen que lidiar y que pueden acabar afectando a los downstream análisis. Ya en diseño experimental, donde se eligen que regiones del gen 16S rRNA hay que secuenciar, existe variabilidad, la cuál, se ve acentuada debido al amplio catálogo de kits de extracción del ADN. A mayores existe mucha variabilidad derivada de la elección de utilizar una determinada plataforma de secuenciación, así como de la estrategia llevada a cabo (Ilumina, PacBIO, Ion Torrent, etc). Finalmente, otro de los puntos clave donde se introduce gran variabilidad es la asignación taxonómica de los datos, donde influyen las pipelines utilizadas (DADA2, Qiime2, MOTHUR, etc), las bases de datos utilizadas (Greengenes, SILVA, RDP, etc) y la estrategia seguida (OTUs o ASVs). Aunque el ML es bueno manejando esta heterogeneidad, inevitablemente se añaden sesgos a los análisis.

Finalmente, la última limitación a destacar es inherente al estudio composicional del microbioma, esto es, al final los estudios que se están realizando se basan en datos procedentes de una instantátena en un momento dado del perfil microbiómico de los pacientes. Intentar reducir la complejidad biológica a un fotograma determinado es simplificar excesivamente procesos dinámicos y multifacéticos de cualquier enfermedad.
Por lo que añadir más puntos temporales (estudios longitudinales) o datos procedentes de otras fuentes ómicas (como metatranscriptoma, metaproteoma o metaboloma) ayudaría a capturar la complejidad e información biológica de manera más realista o representativa.


\section{Futuras líneas de investigación} 

En el futuro, es crucial avanzar hacia una integración efectiva de plataformas analíticas, permitiendo así la unificación de diversos tipos de datos y promoviendo resultados más coherentes y robustos. La combinación de información del microbioma con datos del anfitrión proporcionará una comprensión más integral de las interacciones biológicas y sus implicaciones en la salud. Además, es fundamental investigar cómo el microbioma puede actuar como barrera para la eficacia de ciertos tratamientos, facilitando el desarrollo de estrategias para superar estos obstáculos y mejorar la efectividad terapéutica. Paralelamente, un mayor enfoque en el análisis de resistencia antimicrobiana ayudará a identificar mecanismos clave de resistencia, impulsando el diseño de soluciones innovadoras que optimicen el uso de antimicrobianos y eficienten el tratamiento de infecciones. Asimismo, el desarrollo de modelos longitudinales permitirá predecir cambios dinámicos en el microbioma a lo largo del tiempo, ofreciendo perspectivas valiosas para adaptar tratamientos y mejorar los resultados clínicos.


\section{Cierre de la discusión}


En resumen, las principales contribuciones de la presente Tesis Doctoral se pueden sintetizar en tres ejes fundamentales: la capacidad del microbioma como fuente de biomarcadores, la demostración de la capacidad predictiva y explicable de modelos de ML en el análisis del microbioma, la identificación de  posibles biomarcadores microbianos con potencial de aplicación clínica en el contexto de la DT1/DT2 y el SM, y el desarrollo de una nueva estrategia de subtipificación del microbioma vaginal en el contexto de la VB. Estas aportaciones evidencian la utilidad práctica del enfoque propuesto, haciendo un llamado firme a continuar explorando el microbioma como una herramienta clave para la personalización de intervenciones médicas. Este compendio, en conjunto, sienta las bases para un futuro en el que las tecnologías emergentes y la integración de datos complejos configuren el panorama de la investigación biomédica, aportando beneficios tangibles a la salud pública y a la práctica clínica moderna.


%%%%%%%%%%%%%%%%%%%%%%%%%%%%%%%%%%%%%%%%%%%%%%%          OLD WGS                %%%%%%%%%%%%%%%%%%%%%%%%%%%%%%%%%%%%%%%%%%%%%%%%%%%%%%%%


\subsubsection{Secuenciación del genoma completo}

La secuenciación WGS surge como un enfoque complementario a la secuenciación 16S rRNA, superando ciertas limitaciones de esta última, como los sesgos de \textit{primers} y la resolución limitada. WGS involucra la fragmentación aleatoria del ADN total de una muestra, permitiendo secuenciar ADN bacteriano, viral, fúngico y arqueal sin la necesidad de amplificar un solo gen. Esto facilita un perfil taxonómico más detallado y una caracterización funcional más completa de las comunidades microbianas, identificando genes y rutas metabólicas.

El análisis de lecturas procedentes de WGS generalmente implica dos estrategias complementarias: los enfoques basados en referencias y los enfoques basados en ensamblaje.

\textbf{Enfoques basados en referencias}

En los análisis basados en referencias, las lecturas de secuenciación se procesan mediante algoritmos de mapeo que alinean cada lectura a colecciones de genes y genomas previamente caracterizados en bases de datos curadas. Este enfoque utiliza herramientas de clasificación taxonómica que se valen de estrategias de homología y perfiles de secuencias para asignar una etiqueta taxonómica y funcional a cada lectura o fragmento mapeado. Este enfoque destaca por su rapidez, ya que la comparación se realiza contra un conjunto determinado y bien estudiado de secuencias, lo que reduce el tiempo de cómputo y permite la generación de perfiles detallados en plazos relativamente cortos (de unas horas a varios días, dependiendo del número y la calidad de las muestras, así como de los recursos computacionales disponibles) \cite{bharti2021current}.

Un aspecto importante a considerar es el equilibrio entre rapidez y exhaustividad. Mientras que los métodos basados en referencias destacan por su eficiencia, permitiendo un análisis comparativo rápido entre muestras, su dependencia de bases de datos actualizadas significa que el avance de la investigación en metagenómica requiere, paralelamente, un esfuerzo por parte de la comunidad científica en la incorporación de nuevas secuencias y la corrección de anotaciones erróneas. Esta necesidad de ampliar y actualizar estas bases de datos ha impulsado el desarrollo de técnicas basadas en ensamblaje \cite{bharti2021current}.

\textbf{Enfoques basados en ensamblaje}

En contraste, los enfoques basados en ensamblaje intentan reconstruir secuencias contiguas (\textit{contigs}) de mayor longitud a partir de lecturas solapadas. Este método se fundamenta en algoritmos de ensamblaje (frecuentemente utilizando metodologías de gráficos de De Bruijn \cite{de1948combinatorial} u otros algoritmos de superposición) que integran millones de lecturas para reconstruir fragmentos más largos de genomas completos. La reconstrucción mediante ensamblaje permite, de forma agnóstica a la información previa, obtener secuencias que provienen de organismos hasta entonces no descritos o poco representados en las bases de datos de referencia; facilitando así el descubrimiento de nuevos genes, rutas metabólicas y variantes a nivel de cepa. Esta capacidad de “extender” las lecturas cortas en \textit{contigs} más largos otorga una mayor resolución en la caracterización funcional y taxonómica, ya que se preserva el contexto genómico, lo que resulta fundamental para establecer relaciones entre genes, operones o incluso para identificar elementos móviles \cite{bharti2021current, agustinho2024unveiling}.

No obstante, los métodos basados en ensamblaje también presentan desafíos técnicos y computacionales. En primer lugar, la calidad y la longitud promedio de los \textit{contigs} dependen en gran medida de la calidad intrínseca de las lecturas de secuenciación y de la profundidad de cobertura alcanzada en la muestra. Datos con baja profundidad resultan en ensamblajes fragmentados y con errores, lo cual puede llevar a detectar \textit{contigs} incompletos o a la formación de quimeras. Además, el ensamblaje de muestras metagenómicas, caracterizadas por una alta diversidad y por la existencia de secuencias repetitivas o regiones altamente conservadas entre diferentes organismos, es computacionalmente intensivo, ya que la generación y el procesamiento de grandes grafos de ensamblaje requieren importantes recursos de memoria y de procesamiento \cite{agustinho2024unveiling}.

Al comparar la secuenciación del gen 16S rRNA con WGS, cada enfoque presenta fortalezas y debilidades distintas (ver Figura \ref{fig:sequencing}). Mientras que el 16S rRNA es una opción rentable para el perfilado comunitario a gran escala, la WGS proporciona una visión holística y funcional de la comunidad microbiana. La complejidad, el costo y los requisitos computacionales de la WGS la hacen menos accesible para estudios de rutina, pero son invaluables para la caracterización profunda del microbioma.

%% Justificación 16S Tesis Doctoral %%

Esta Tesis Doctoral se ha centrado principalmente en el uso de datos 16S rRNA debido a varias razones. Primero, la secuenciación 16S rRNA es más rentable y accesible, permitiendo realizar estudios a gran escala sin un coste prohibitivo. Además, requiere menos capacidad computacional, facilitando un análisis más eficiente y rápido. Esto es especialmente importante en la identificación preliminar de biomarcadores, donde se busca caracterizar rápidamente comunidades microbianas en diversas muestras. Asimismo, es particularmente valiosa para estudios que priorizan la reproducibilidad; su metodología estandarizada permite comparaciones consistentes entre estudios y a lo largo del tiempo, algo esencial para validar posibles biomarcadores. Finalmente, la abundancia de conjuntos de datos disponibles de 16S rRNA es mucho más elevada que de WGS. Esto habilita el desarrollo de trabajos más complejos (como validaciones externas o estudios de metacohortes) para el estudio de diferentes enfermedades. Aunque ofrece menos resolución funcional que WGS, su pe


%%%%%%%%%%%%%%%%%%%%%%%%%%%%%%%%%%%%%%%%%%%%%%%        OLD  Objetivos                %%%%%%%%%%%%%%%%%%%%%%%%%%%%%%%%%%%%%%%%%%%%%%%%%%%%%%%%


\chapterintro{
El desarrollo de esta Tesis Doctoral está motivado por la necesidad de avanzar en el campo de la bioinformática aplicada a la salud humana. En el contexto actual, la identificación de biomarcadores microbianos es fundamental para mejorar el diagnóstico y pronóstico de diversas enfermedades. Este capítulo detalla los objetivos que guían este trabajo de investigación, fundamentando la importancia de cada uno en el ámbito científico y médico.}


El objetivo principal de esta Tesis Doctoral es el \textbf{desarrollo y aplicación de metodologías bioinformáticas para la identificación de biomarcadores microbianos asociados a enfermedades humanas}. Esto se aborda con futuros fines diagnósticos, utilizando el análisis de datos de microbioma obtenidos mediante secuenciación masiva y técnicas de Machine Learning.


Por consiguiente, a partir del objetivo principal se plantean los siguientes objetivos específicos:
\begin{itemize}
    \item \textbf{Establecer e implementar estrategias bioinformáticas y modelos computacionales que permitan identificar patrones microbianos específicos de enfermedades}. Este objetivo se centra en la creación de metodologías y flujos de trabajo para el análisis exhaustivo de datos del microbioma humano. Se propone desarrollar herramientas que faciliten el procesamiento desde datos crudos hasta la identificación de potenciales biomarcadores. Estas herramientas abarcan tanto análisis de abundancias microbianas como técnicas de ML para detectar y caracterizar patrones microbianos informativos dentro de una determinada patología. Con la implementación de estas estrategias se busca una interpretación más precisa de los datos con el fin de comprender mejor el vínculo entre el microbioma y la homeostasis del individuo.
    
    \item \textbf{Explorar el uso del microbioma humano como fuente de biomarcadores mediante el análisis de datos metagenómicos en distintas patologías humanas}. Este objetivo busca investigar cómo las comunidades microbianas pueden reflejar cambios relacionados con diversas enfermedades. Se busca entender las complejidades y dinámicas de las comunidades microbianas, evaluando si sus cambios pueden servir como indicadores de alteraciones patológicas. Además, se estudia la viabilidad de utilizar estos cambios microbianos para estratificar o mejorar la clasificación de pacientes en estudios clínicos, con el fin de contribuir al avance en métodos de detección y tratamiento basados en el microbioma.

    \item \textbf{Validar los biomarcadores microbianos identificados mediante un análisis comparativo en cohortes clínicas independientes}. Este objetivo se centra en evaluar la consistencia y robustez de los biomarcadores identificados inicialmente en una cohorte principal mediante comparaciones in silico con datos de cohortes externas. Esta validación externa garantiza que los biomarcadores mantengan su relevancia y potencial clínico en diferentes contextos, fortaleciendo su aplicabilidad en estudios futuros.
    
\end{itemize}

%%%%%%%%%%%%%%%%%%%%%%%%%%%%%%%%%%%%%%%%%%%%%%%        OLD  Discusion                %%%%%%%%%%%%%%%%%%%%%%%%%%%%%%%%%%%%%%%%%%%%%%%%%%%%%%%%


\section{Introducción general de la discusión}

\textbf{Breve recapitulación del objetivo principal de la tesis}

El objetivo principal de esta Tesis Doctoral ha sido explorar el potencial del microbioma humano como fuente de futuros biomarcadores clínicos, utilizando técnicas computacionales. A través de un análisis detallado de datos metagenómicos, se busca identificar patrones específicos que puedan asociarse con diversas condiciones de salud. Asimismo, se busca crear modelos predictivos a partir de estos patrones específicos que permitan clasificar a pacientes de determinadas patologías. Este enfoque no solo pretende ampliar el conocimiento sobre la relación entre el microbioma y estas patologías, sino también establecer posibles vías alternativas para futuras investigaciones en medicina.


\textbf{Relevancia del tema}

El estudio del microbioma humano ha cobrado gran relevancia en los últimos años debido a su potencial como fuente de biomarcadores para enfermedades complejas. Como se ha expuesto en la presente Tesis Doctoral, cada vez hay más evidencias que defienden que el microbioma modula funciones esenciales del sistema inmune, el metabolismo y la respuesta inflamatoria. Esto se refleja en investigaciones que han asociado alteraciones en el microbioma con patologías como la DT2, algunos tipos de cáncer y trastornos autoinmunitarios \cite{schirmer2016linking, fan2021gut, bakir2022inflammatory}. Por lo que, los biomarcadores microbianos pueden ofrecer una ventana para la detección temprana y un entendimiento más profundo de estas enfermedades. Además, el estudio en profundidad de estos biomarcadores y su comparación con los ya existentes en el estado del arte puede llevar a mejoras en las estrategias diagnósticas o terapéuticas. 

\textbf{Justificación metodológica general}

El enfoque metodológico adoptado en la presente Tesis Doctoral se fundamenta en la integración de técnicas estadísticas de análisis de abundancia diferencial, técnicas de selección de características y técnicas de ML para abordar las particularidades y limitaciones inherentes al análisis del microbioma humano. Por una parte, los métodos de abundancia diferencial específicos para microbiomas, al utilizar transformaciones logarítmicas y modelos estadísticos adaptados a la sobredispersión e inflación de ceros, aseguran que las inferencias realizadas sean consistentes con las propiedades intrínsecas de los datos. Por otra parte, la utilización de estrategias de selección de características posibilita la reducción drástica de la dimensionalidad del conjunto de datos sin sacrificar la información crítica, facilitando la extracción de biomarcadores que exhiben alta capacidad discriminativa y relevancia biológica. Finalmente, la estrategia de ML seguida permite construir modelos predictivos capaces de manejar la elevada variabilidad en las comunidades microbianas entre individuos, posibilitando así la identificación de firmas microbianas que pueden predecir estados patológicos con elevado poder de clasificación y generalización.


\section{Evaluación crítica transversal de los resultados}

\textbf{1. Enfermedades}

En primera instancia, se decidió estudiar las enfermedades presentadas a lo largo de la Tesis Doctoral debido a su conocida y fuerte conexión con algún componente microbiano. En el caso de las enfermedades metabólicas, estas tienen una relación directa con el tracto intestinal y su correcto funcionamiento, un ambiente que presenta una gran carga microbiana. En cuanto a la VB, se sabe que la patogénesis de la enfermedad tiene un gran porcentaje microbiano. Por ello, analizar estas enfermedades desde esta perspectiva podría ser sumamente relevante para intentar comprender mejor sus mecanismos subyacentes.

Al estudiar estas enfermedades \cite{paper1, paper2, paper3}, es fundamental considerar la naturaleza de las muestras utilizadas. Las muestras fecales, al ser deslocalizadas, reflejan una visión general del microbioma intestinal a lo largo de todo el tracto digestivo. Esto es crucial para enfermedades sistémicas y complejas como el SM o DT1/2, donde el impacto puede deberse a interacciones a gran escala dentro del sistema digestivo. Sin embargo, su carácter heterogéneo puede complicar la identificación de microorganismos específicos que puedan ser responsables de un efecto observado. En contraste, las muestras localizadas, como los hisopos vaginales, ofrecen un análisis más directo de un ecosistema microbiano específico y menos heterogéneo. Esto es particularmente útil para estudios donde el microambiente es un factor clave, como en la VB. Dado que el microbioma vaginal es menos diverso, las alteraciones en su composición son más fáciles de asociar directamente con alteraciones específicas en la salud vaginal.

Por un lado, es relevante destacar el papel de la VB como modelo de enfermedad con un origen claramente asociado a la disbiosis microbiana. En este caso, la alteración del equilibrio en la comunidad vaginal, caracterizada por una disminución de \textit{Lactobacillus} y un aumento de bacterias anaerobias, no solo se asocia de forma consistente con la condición clínica, sino que cumple un rol causal bien establecido en su fisiopatología. Por otro lado, la identificación de firmas microbianas robustas en cohortes de DT1 y DT2 refuerza la hipótesis de que el microbioma refleja (y modula) estados fisiopatológicos sistémicos. No obstante, dado que en estas condiciones la relación microbioma-enfermedad es más compleja, multifactorial y posiblemente bidireccional, los resultados deben ser interpretados con cautela.

En conjunto, la selección de enfermedades con diferentes grados de implicación microbiana ha permitido evaluar la versatilidad de los enfoques computacionales empleados, pero también pone de manifiesto la necesidad de futuras investigaciones para validar y generalizar estos hallazgos en otras condiciones clínicas.

\textbf{2. Datos}

Respecto a la elección de utilizar datos procedentes de 16S rRNA para estudiar estas enfermedades, se deben considerar varias razones. Al ser una metodología más económica, el número de trabajos científicos (y de cohortes) con los que comparar nuestro enfoque y resultados es mucho mayor que si se hubiera optado por usar datos de WGS. En relación con esto, en la búsqueda de potenciales biomarcadores clínicos, la reproducibilidad de los resultados obtenidos resulta más accesible económicamente al emplear esta metodología en lugar de WGS.

Ligado a esto, para inferir los taxones microbianos procedentes de 16S rRNA, se decidió utilizar ASVs por ser el enfoque con mayor resolución posible, ya que podría facilitar la búsqueda de nuevos biomarcadores potenciales. Además, la corrección de secuencias y la mayor reproducibilidad que ofrecen las ASVs sobre las OTUs también fueron factores importantes que se consideraron. Finalmente, la tendencia de la comunidad científica a favor del uso de ASVs sobre OTUs influyó en la toma de decisiones \cite{callahan2017exact}.

Más específicamente, en lo que respecta a la elección del método de inferencia de ASVs, se seleccionó DADA2, dado que se buscaba un método que infiriera el mayor número de ASVs posibles sin perder sensibilidad al representar fielmente las comunidades microbianas reales. En relación con esto, numerosos artículos en los que se comparan diferentes métodos de inferencia señalan a DADA2 como el método con mayor sensibilidad y el que mejor modela los errores procedentes de la secuenciación. Otro punto importante para decidir utilizar este método es que es el único independiente del conjunto de datos; es decir, los resultados de la inferencia de secuencias para una muestra no varían en función de las muestras que se procesan con ella. Esto representa un punto clave para la reproducibilidad de experimentos \cite{nearing2018denoising, caruso2019performance, prodan2020comparing}.

\textbf{3. Composicionalidad}

Respecto a la selección del enfoque de procesamiento de los conteos taxonómicos, la mayor preocupación era intentar paliar la gran variabilidad en la profundidad de secuenciación entre muestras y representar fielmente las proporciones de los organismos microbianos. Se observó que métodos como la rarefacción (el más utilizado hasta hace unos años) pueden simplificar en exceso las comparaciones al eliminar datos potencialmente relevantes, reduciendo la potencia estadística del análisis y subestimando la diversidad real. Además, imponen un sesgo debido al submuestreo aleatorio, afectando especialmente la detección de taxones poco abundantes. En contraposición, se observó que el enfoque composicional conserva la totalidad de los datos sin recurrir al submuestreo, respeta la naturaleza composicional del microbioma y mejora la capacidad para detectar diferencias reales entre las comunidades microbianas. Debido a las ventajas que aporta este último enfoque, se decidió escogerlo para el procesamiento de los conteos taxonómicos, así como el uso de métodos de abundancia diferencial que integren este enfoque \cite{hong2022rarefy, gloor2016compositional, gloor2017microbiome, zhou2022linda}.

Actualmente, en la literatura todavía se siguen utilizando métodos de normalización para datos procedentes de 16S rRNA que no tienen en cuenta la composicionalidad de los datos. Numerosos trabajos continúan abogando por la rarefacción a un número fijo de lecturas; tenemos ejemplos tanto en diabetes \cite{lambeth2015composition, dash2023functional, siptroth2023variation} como en BV \cite{chen2018vaginal, onywera2021predictive}, lo que introduce un sesgo, en mayor o menor medida, en los resultados obtenidos. Gracias a dichas transformaciones y al uso de LinDA en el primer trabajo \cite{paper1}, se extrajeron patrones significativos. En el caso de los otros trabajos \cite{paper2, paper3}, se comparó el rendimiento aplicando estas transformaciones frente a los datos crudos y se observó que el rendimiento se reducía al no aplicar dichas transformaciones. Además, en este último trabajo, aplicar estas transformaciones ayuda a reducir la variabilidad entre los diferentes enfoques utilizados (OTUs vs ASVs) a lo largo de las distintas cohortes.

\textbf{4. Enfoque Data-driven frente a Knowledge-driven: ML}

En el caso de seleccionar el tipo de enfoque para analizar los datos, esta selección depende de los objetivos de la presente Tesis Doctoral. Un enfoque basado en datos ofrece mucha más flexibilidad y libertad para descubrir nuevos potenciales biomarcadores que un enfoque basado en conocimiento. Al no estar limitado por las restricciones y supuestos de un enfoque basado en conocimiento, se pueden explorar los datos sin hipótesis predefinidas, lo que permite el hallazgo de patrones y relaciones desconocidas.

El enfoque dirigido por conocimiento depende de hipótesis preexistentes, vías biológicas conocidas o bases de datos curadas. Aunque proporciona interpretabilidad biológica y es un excelente complemento para confirmar hallazgos, limita la capacidad de descubrimiento a lo ya conocido \cite{dakal2025advanced}. El enfoque dirigido por datos, adoptado aquí principalmente a través de técnicas de ML, opera sin depender de suposiciones biológicas. Esto es crucial en la búsqueda exploratoria de nuevos biomarcadores, ya que permite analizar datos brutos y descubrir patrones ocultos o no lineales \cite{hemmati2024using}.

En el estado del arte actual, se utilizan muchos enfoques basados en el conocimiento en el campo del microbioma \cite{chong2017computational, bauer2018network, luo2024progress}. Sin embargo, a lo largo de los tres trabajos \cite{paper1, paper2, paper3}, al analizar datos brutos en lugar de un subconjunto basado en conocimiento previo, hemos descubierto tanto taxones ya mencionados en la literatura como otros poco documentados o desconocidos en la patología de estudio. Esto abre la puerta a investigar el papel de estos taxones en el desarrollo de enfermedades y su potencial como biomarcadores.

A la hora de escoger una metodología guiada por este tipo de enfoque, se selecciona el ML, ya que se deseaba utilizar una metodología que manejase, en cierta medida, la alta variabilidad, dispersión y complejidad del microbioma, y que además tuviese explicabilidad. Se descartaron otros enfoques como el \textit{Deep Learning}, ya que es un enfoque de caja negra que no ofrece explicabilidad y no se contaba con un número adecuado de muestras, lo que probablemente iba a propiciar un sobreajuste de las redes entrenadas.


\textbf{5. Dimensionalidad}

Para afrontar el desafío de la alta dimensionalidad en los datos del microbioma, se han implementado filtrados rigurosos (conteos totales y prevalencia de taxones) al extraer la tabla de taxones y se han aplicado técnicas de selección de características que permiten reducir el número de taxones considerados en la construcción de modelos de ML, las cuales han demostrado ser efectivas en este contexto \cite{segata2011metagenomic}. Este enfoque optimiza los modelos al disminuir el ruido y las variables redundantes, mejorando también su interpretabilidad y rendimiento.

Actualmente, el problema de la dimensionalidad está presente en el estudio de datos procedentes del microbioma. Numerosos artículos en el estado del arte que utilizan enfoques basados en ML suelen realizar algún tipo de selección de características, ya sea a través de algoritmos que realizan su propia selección (como RF o glmnet) o aplicando algún método de selección y ordenación de características según su relevancia para un problema específico \cite{beck2015machine, biassoni2020gut, bakir2021discovering, bakir2022inflammatory, chavarria2025using}. Esta selección de características desempeña un papel importante en el segundo trabajo de esta Tesis Doctoral \cite{paper2}. Este artículo emplea la selección de características para identificar las especies más relevantes capaces de diferenciar las muestras entre sano y DT1. Los resultados muestran un aumento significativo en el rendimiento de los modelos al seleccionar un subgrupo de especies del número total disponible. La estrategia seguida fue utilizar una técnica de selección de características de filtrado, lo que permite medir la relevancia de cada una de las especies con las clases de interés, identificando así cuántas y qué especies utilizar para construir el modelo predictivo.


\textbf{6. Variabilidad}

La variabilidad entre muestras constituye otro obstáculo significativo en el análisis del microbioma, un reto que se ha abordado mediante el uso de técnicas avanzadas de validación. En el caso del aprendizaje no supervisado, el enfoque CCP ha permitido la construcción de subgrupos robustos y generalizables. Gracias al submuestreo, a la evaluación de los algoritmos de \textit{clustering} y a la evaluación consensuada de los \textit{clusters} producidos, se asegura que el modelo final de \textit{clustering} sea robusto y no presente sobreentrenamiento en un subgrupo de los datos originales. En el aprendizaje supervisado, la implementación de una NCV garantiza un modelo robusto y generalizable que contemple la variabilidad inherente a las muestras biológicas \cite{varma2006bias}. Al segmentar el conjunto de datos en múltiples subconjuntos de entrenamiento y prueba, se logra una estimación fiable del rendimiento del modelo y se minimiza el riesgo de sobreajuste.

Esto se refleja en los trabajos \cite{paper2, paper3}, donde, gracias a esta validación, se ha manejado la diversidad en los datos de entrada, permitiendo la construcción de modelos con poca dispersión en su rendimiento (independientemente del pool de entradas). Esto se traduce en un modelo con altas capacidades para clasificar pacientes entre sanos y DT1 en el trabajo \cite{paper2}, y en un modelo en el trabajo \cite{paper3} que es capaz de predecir los VCS, incluso cuando los datos de entrada provienen de fuentes heterogéneas.

\textbf{7. Dispersión}

Para mitigar la alta dispersión presente en los datos del microbioma, se ha promovido el uso de algoritmos de ML que han demostrado ser altamente eficaces en este contexto \cite{goodswen2021machine, Lin2023MicrobiomeObesitySignatures}. Los RF, al generar y combinar múltiples árboles de decisión, permiten reducir la varianza y mejorar la estabilidad del modelo, siendo efectivos frente a la dispersión. Asimismo, algoritmos como XGBoost han demostrado ser robustos frente a \textit{outliers} y datos ruidosos debido a su capacidad de optimización iterativa con regularización, lo que facilita el manejo de datos dispersos. Estos algoritmos no solo gestionan la dispersión intrínseca de los datos microbianos, sino que también mejoran la precisión y la generalización de los modelos predictivos.

El uso de estos algoritmos ha sido clave en los trabajos \cite{paper2, paper3}. Gracias a la implementación de algoritmos que se regularizan internamente, se ha evitado que taxones sin importancia, o que puedan confundir a los modelos, afecten negativamente mediante la penalización correspondiente.

\textbf{8. OTUs vs ASVs?}

Uno de los puntos de partida de la investigación sobre el microbioma es la elección del enfoque para la inferencia de secuencias. En esta Tesis Doctoral, se ha explorado la utilidad de las OTUs y las ASVs para el análisis del microbioma \cite{paper1, paper2, paper3}.

En el primer trabajo \cite{paper1}, se utilizaron ASVs para analizar la abundancia diferencial entre controles y DT2. Su alta granularidad permitió detectar señales sutiles que la agrupación en OTUs podría diluir. Por otro lado, en el segundo trabajo \cite{paper2}, el análisis con OTUs identificó una firma metagenómica en DT1. Aunque la agrupación puede fusionar secuencias e incrementar la riqueza, las OTUs capturan patrones robustos a niveles superiores \cite{liu2023independent}. En este contexto, fueron suficientes para un modelo de clasificación altamente predictivo. Además, en el último trabajo \cite{paper3}, un modelo de \textit{clustering} basado en OTUs se validó con datos procesados con ASVs, mostrando que los patrones comunitarios principales se mantienen entre ambos enfoques. Esto está en línea con evidencias que indican que, pese a diferencias en riqueza o asignación fina, la beta diversidad, la topología del \textit{clustering} y las señales globales tienden a ser comparables \cite{chiarello2022ranking}.

En conjunto, las discusiones en estos tres artículos destacan que, aunque los enfoques basados en OTUs y ASVs difieren en sus fundamentos teóricos y resultados prácticos, ambos pueden generar resultados robustos y biológicamente significativos cuando se aplican cuidadosamente protocolos de preprocesamiento y filtrado. La mayor resolución de los ASVs resulta especialmente valiosa para detectar diferencias sutiles en la abundancia entre grupos de estudio, especialmente cuando se buscan especies singulares o incluso cepas, como se evidencia en \cite{paper1}. Por otro lado, la robustez y la simplicidad del agrupamiento en OTUs han demostrado ser suficientes para extraer firmas metagenómicas predictivas para la clasificación de enfermedades, como se observa en \cite{paper2}, y para capturar la agrupación a nivel comunitario en \cite{paper3}. En última instancia, estos hallazgos refuerzan la idea de que la elección del método debe alinearse con los objetivos del estudio.


\textbf{9. Modelos lineales vs modelos no lineales} 

En el segundo trabajo \cite{paper2}, la tarea de predecir pacientes DT1 utilizando datos del microbioma intestinal presenta un gran desafío, ya que el ecosistema intestinal es altamente complejo, con interacciones microbianas no lineales y multifactoriales por naturaleza. En este contexto, el microbioma intestinal refleja una amplia gama de taxones en interacción cuyas relaciones son moduladas por la colonización del nicho, factores ambientales abióticos y respuestas inmunitarias del huésped. Los algoritmos de ML no lineales, como RF, están inherentemente equipados para capturar estas complejas estructuras de dependencia, ya que operan mediante la construcción de un conjunto de árboles de decisión que dividen jerárquicamente el espacio de características. Este enfoque en conjunto permite a los modelos RF identificar interacciones y asociaciones de orden superior entre características microbianas que pueden ser críticas para discriminar entre estados de enfermedad y condiciones saludables.

Por otro lado, el tercer trabajo \cite{paper3} se centra en la estratificación de pacientes utilizando perfiles del microbioma vaginal y emplea un modelo lineal, como glmnet. Como se ha comentado, el microbioma vaginal se caracteriza por ser mucho menos heterogéneo, presentar una diversidad reducida y donde las comunidades microbianas tienden a estar dominadas por especies de \textit{Lactobacillus}. En tales casos, los límites que separan los perfiles comunitarios asociados con diferentes grupos de pacientes parecen ser más linealmente separables, lo que permite que modelos lineales como glmnet capturen efectivamente estas diferencias con el beneficio adicional de la interpretabilidad del modelo. La claridad en la delimitación de la comunidad microbiana en el nicho vaginal con frecuencia conduce a una situación en la que la suposición más simple de efectos aditivos es suficiente para la estratificación de pacientes, reduciendo así la necesidad de un modelado no lineal más complejo.

En resumen, el rendimiento superior de algoritmos no lineales como el RF en el contexto de la predicción de DT1 basada en el microbioma intestinal \cite{paper2} se debe principalmente a la necesidad de modelar interacciones complejas del ecosistema microbiano intestinal, las cuales se ven agravadas por una alta variabilidad interindividual y las interacciones anfitrión-huésped. En contraste, el uso efectivo de un modelo lineal, como glmnet, en la estratificación de pacientes en función de los perfiles del microbioma vaginal \cite{paper3} refleja la naturaleza relativamente estructurada y linealmente separable de las comunidades microbianas vaginales, donde se pueden delimitar características clave con precisión, aprovechando los beneficios de la interpretabilidad y la escasez del modelo. Las diferencias en el contexto de la enfermedad y la heterogeneidad de los nichos microbianos son, por lo tanto, fundamentales para explicar por qué los métodos no lineales superan en un escenario, mientras que los modelos lineales son más efectivos en otro.

\textbf{10. Microbioma longitudinal}

Como se ha mencionado, el microbioma humano es un ecosistema altamente dinámico que exhibe variabilidad temporal, tanto a nivel composicional como funcional, una característica reconocida en la literatura que plantea desafíos para la clasificación confiable de enfermedades utilizando datos del microbioma \cite{aldars2021systematic}. Los estudios \cite{paper2, paper3}, al igual que otros \cite{eck2017robust}, han demostrado que, incluso utilizando solo una muestra por paciente para la clasificación de enfermedades, se pueden lograr resultados robustos. Esto sugiere que la señal microbiana relacionada con la patología puede ser lo suficientemente resistente como para emerger a pesar de las fluctuaciones temporales subyacentes. Es importante destacar que, aunque el microbioma experimenta cambios pronunciados a lo largo de horas, días y semanas, estudios han revelado que las características centrales que proporcionan utilidad diagnóstica son, de hecho, más estables de lo que se podría esperar \cite{mehta2018stability}.

Los resultados obtenidos, a pesar del comportamiento variable del microbioma, indican que los modelos utilizados probablemente están aprovechando características que representan ya sea una respuesta promedio o aspectos altamente conservados del microbioma que son menos sensibles a las perturbaciones temporales. En esencia, aunque una única instantánea puede no captar todo el espectro de fluctuaciones temporales, parece encapsular un estado basal suficientemente representativo de las firmas microbianas relevantes para la enfermedad.

Esta observación concuerda con la evidencia reportada en varios estudios, que sugieren que la robustez de la clasificación de enfermedades no depende únicamente de capturar cada fluctuación transitoria, sino más bien de detectar los marcadores microbianos persistentes y dominantes que diferencian la salud de la enfermedad.


\section{Evaluación crítica de los artículos}

\textbf{Primer trabajo}

En el primer trabajo \cite{paper1}, los resultados en pacientes con SM revelaron un aumento significativo de la relación Firmicutes/Bacteroidetes. Este hallazgo está en concordancia directa con la hipótesis clásica que asocia el aumento de dicha relación con la obesidad y los fenotipos metabólicos \cite{santos2019role, baghi2022evaluation}. Se observó una mayor prevalencia del género \textit{Blautia} en pacientes con SM. Aunque algunos estudios sugieren un aumento en la abundancia de \textit{Blautia} en individuos con SM, los mecanismos subyacentes de esta asociación siguen sin estar claros. Investigaciones previas han indicado que una mayor presencia de \textit{Blautia} puede contribuir a la disfunción metabólica a través de vías como la inflamación y la disrupción de la barrera intestinal \cite{Liu2021}. Además, se ha identificado una correlación positiva entre la abundancia de \textit{Blautia} y la severidad de SM en pacientes con enfermedad del hígado graso no alcohólico, una condición estrechamente relacionada con SM \citep{shen2017gut, Wang2016}. El enfoque seguido permitió la identificación de nuevos taxones que podrían desempeñar un papel protector en el contexto del SM. En particular, se relacionó la especie \textit{Parabacteroides merdae} con los controles sanos. Aunque no ha sido extensamente estudiada en humanos, trabajos recientes en modelos murinos sugieren que esta especie podría desempeñar un rol beneficioso a través del catabolismo de aminoácidos de cadena ramificada, cuya acumulación se ha relacionado con un aumento del riesgo metabólico \cite{qiao2022gut}. Además, estudios previos han reportado una reducción de ciertas especies del género \textit{Parabacteroides} en individuos con SM \cite{cui2022roles}.

La progresión a DT2 mostró un aumento en la relación Bacteroidetes/Firmicutes y una notable disminución del filo Actinobacteria. Esta disminución implica la pérdida de géneros productores de butirato, como \textit{Bifidobacterium}, un taxón ampliamente estudiado en el contexto de la DT2 \cite{larsen2010gut, wu2010molecular, le2013alterations}. La pérdida de esta funcionalidad protectora se vio reforzada por la reducción en la proporción de dos géneros poco conocidos en este contexto: \textit{Lachnospiraceae GCA-900066575} y \textit{Defluviitaleaceae UCG-011}. Al igual que otros miembros del orden \textit{Clostridiales}, estos podrían participar en la fermentación de fibras y la producción de AGCCs, que, como ya se ha visto, tienen efectos antiinflamatorios y reguladores sobre la función de la barrera intestinal. Además, estudios muy recientes realizados en ratones señalan que \textit{GCA-900066575} se relaciona negativamente con niveles de glucosa en sangre en ayunas y la resistencia a la insulina \cite{liu2023vitamin}. Simultáneamente, se observó un incremento en los taxones asociados con la patogénesis, como \textit{Dorea}, \textit{Prevotella} y \textit{Dialister invisus}. En el caso específico de \textit{Dorea longicatena}, se ha encontrado que esta especie está enriquecida en pacientes con DT2. Se ha observado una correlación clínica significativa entre su presencia y el desarrollo de la enfermedad \citep{qin2014alterations}. En cuanto a \textit{Prevotella}, varios estudios han establecido una asociación entre un aumento de esta bacteria y la intolerancia a la glucosa, así como la diabetes gestacional \citep{zhang2013human, egshatyan2016gut, hasain2020gut}. Por otro lado, la especie \textit{Dialister invisus} se ha vinculado con un aumento de la permeabilidad intestinal y la diabetes \citep{maffeis2016association, zheng2018gut}.

\textbf{Segundo trabajo}

El segundo trabajo \cite{paper2} muestra que la profundidad taxonómica del estudio es un factor determinante para capturar la señal asociada a la T1D. Se observa que el rendimiento de los tres algoritmos utilizados (RF, SVM y glmnet) mejora a través de seis niveles taxonómicos, desde filo hasta especie, con mejoras del 30 \% entre ambos niveles. Este patrón sugiere que la asociación con la enfermedad se basa en la dinámica o interacción entre especies, en lugar de en desequilibrios de niveles superiores. En cuanto a la selección de características a nivel de especie, el método de Wilcoxon permite identificar firmas compactas para cada algoritmo (RF = 25, SVM = 28, glmnet = 15). Al entrenar los algoritmos con estas firmas, se evidencia que RF supera significativamente a SVM y glmnet (p-valor = 0.0079), manteniendo un AUC muy estable ($<$ 0.97). Por consiguiente, la firma de estas 25 especies se establece como el equilibrio óptimo entre rendimiento e interpretabilidad biológica.

La firma metagenómica identificada se alinea con asociaciones previamente descritas a nivel de especie. \textit{Bacteroides dorei} se ha vinculado con la seroconversión y con etapas preclínicas de autoinmunidad \cite{davis2014bacteroides}. \textit{Prevotella copri} y \textit{Bacteroides vulgatus} se han asociado con perfiles metabólicos de resistencia a la insulina \cite{pedersen2016human}. En cuanto a \textit{Bifidobacterium animalis}, se han observado efectos favorables en ensayos clínicos (principalmente en DT2) sobre control glucémico y parámetros cardiometabólicos \cite{stenman2016probiotic, khiaolaongam2025bifidobacterium}. En el eje intestino-inmunidad, \textit{Bacteroides thetaiotaomicron} se ha asociado con fenotipos de enfermedad celíaca y presenta una capacidad inmunomoduladora relevante \cite{sanchez2010intestinal}. Por último, la dinámica de \textit{Escherichia coli} productora de amiloide se ha relacionado con vías prodiabéticas y con la seroconversión en cohortes pediátricas \cite{tetz2019type}.

Los modelos entrenados con esta firma presentan un gran rendimiento (0.94 de \textit{accuracy}) tanto al clasificar entre individuos sanos y pacientes de DT1, como al añadir una clase intermedia de seroconvertidos. Esto sugiere una fuerte asociación entre estas especies bacterianas y la susceptibilidad a DT1. En este contexto, la mayoría de los estudios describen alteraciones composicionales y funcionales asociadas a la seroconversión y DT1 \cite{vatanen2018human, kostic2015dynamics}. De forma similar, otros estudios poblacionales \cite{belteky2023infant} han aplicado algoritmos de ML (RF) para la clasificación de taxones en búsqueda de biomarcadores; sin embargo, no se detalla su desempeño predictivo para poder comparar con este trabajo.

Desde la publicación de este segundo trabajo, en los últimos años se han publicado diferentes estudios con nuevas cohortes de infantes (como PRJNA875929 \cite{belteky2023infant} y PRJNA702261 \cite{del2022gut}), dando pie a realizar estudios de metacohortes con enfoques similares al presentado en esta Tesis Doctoral, como es el caso de \cite{liu2024predicting}.


\textbf{Tercer trabajo}


En el tercer trabajo \cite{paper3}, los subgrupos obtenidos mediante \textit{VIBES} (VCS I–IV) mostraron una alta estabilidad en la cohorte de descubrimiento, con un ARI medio de 0,914. Esto sugiere que la estratificación en cuatro grupos está determinada por la estructura intrínseca de los datos, y no por el tamaño muestral. Además, la distribución de los subtipos fue consistente en cinco cohortes externas. En la validación externa del modelo, la BA fue superior a 0,91 en todas las cohortes, y el \textit{kappa} fue mayor a 0,83. En conjunto, estos resultados respaldan la portabilidad del modelo y su utilidad para asignar subtipos en datos no vistos.

En la caracterización de los subtipos obtenidos, las betas del modelo cuantifican la contribución relativa de cada especie en el clasificador para determinar la pertenencia a un subtipo. En VCS-I, \textit{Lactobacillus crispatus} presenta una beta dominante, consistente con un ecosistema saludable caracterizado por una abundante producción de ácido láctico, un pH bajo y la exclusión competitiva de patobiontes \cite{ravel2011vaginal, srinivasan2012bacterial}. Por otro lado, en VCS-II destacan especies como \textit{Lactobacillus iners}, \textit{Aerococcus christensenii}, \textit{Campylobacter ureolyticus} y \textit{Staphylococcus haemolyticus}, siendo \textit{Lactobacillus iners} un \textit{lactobacilo} controvertido que se detecta tanto en estados de salud como de disbiosis; su presencia sugiere un estado de transición \cite{macklaim2011crossroads, holm2023lactobacillus}.

En VCS-III, las betas relevantes incluyen \textit{Lactobacillus iners}, \textit{Prevotella bivia}, \textit{Campylobacter ureolyticus} y \textit{Finegoldia magna}. Las especies del género \textit{Prevotella} pueden elevar el pH mediante la producción de aminas biógenas (putrescina, cadaverina, trimetilamina) y amonio, favoreciendo el crecimiento de \textit{Gardnerella} y promoviendo la transición hacia la disbiosis \cite{miller2016lactobacilli, yeoman2013multi}. Finalmente, VCS-IV es el subtipo más heterogéneo, evidenciando altas betas para \textit{Atopobium vaginae}, \textit{Sneathia sanguinegens}, \textit{Mycoplasma hominis} y \textit{Campylobacter ureolyticus}, y en menor medida para \textit{Lactobacillus gasseri}. La especie \textit{Atopobium vaginae} está asociada a \textit{biofilms}, elevado pH y recurrencias debido a su menor sensibilidad al metronidazol \cite{swidsinski2008adherent, bradshaw2006association}; mientras que \textit{Sneathia} expresa factores de virulencia con potencial inflamatorio \cite{fredricks2005molecular, harwich2012genomic}.

Metodológicamente, \textit{VIBES} y \textit{VALENCIA} se basan en paradigmas distintos. \textit{VALENCIA} utiliza el algoritmo \textit{K-nearest neighbors} y hasta 199 taxones, desde clase hasta especie, para asignar sus CSTs en 5 clases y subclases \cite{france2020valencia}. Por otro lado, \textit{VIBES} aplica un modelo \textit{Elastic-net} con un conjunto reducido de 20 especies, definiendo 4 subtipos. Esta diferencia en el algoritmo, en el espacio de taxones y en la estructura de clases permite captar señales no redundantes. Aunque los perfiles sanos/asintomáticos de \textit{VALENCIA} tienden a mapear con VCS-I, varios CSTs (III, IV-A, V) no se alinean de manera directa con VCS-II/III, que presentan mayor heterogeneidad; esto aporta una granularidad complementaria para la estratificación \cite{ravel2011vaginal}. En la cohorte longitudinal tratada con metronidazol, el modelo combinado (\textit{VALENCIA} + \textit{VIBES}) mejora la predicción de respuesta en comparación con cada enfoque por separado. Además. los no respondedores se concentran en VCS-IV, y en menor medida en VCS-III, sugiriendo que los subtipos de \textit{VIBES} añaden información clínica independiente.

\section{Limitaciones del estudio}

Un análisis crítico de la presente Tesis Doctoral requiere reconocer sus limitaciones, esenciales para contextualizar los hallazgos y orientar futuras investigaciones. Estas limitaciones se identifican en tres niveles: primero, en cada uno de los artículos; segundo, en las diversas fuentes de variabilidad que pueden influir en los datos; y tercero, en el estudio intrínseco del microbioma.

\textbf{Limitaciones de los trabajos}

En el primer trabajo \cite{paper1}, el poder estadístico del análisis puede verse afectado por el tamaño muestral de la cohorte (n=79). Además, la distribución asimétrica de sujetos entre las clases de estudio (individuos sanos, prediabéticos y con diabetes tipo 1) es un factor a considerar, ya que podría influir en la robustez de la metodología para detectar diferencias significativas. Adicionalmente, al tratarse de una cohorte de una región geográfica específica, podría no reflejar la heterogeneidad de poblaciones más amplias, limitando así la capacidad de generalización de los resultados.

Respecto al segundo trabajo \cite{paper2}, su principal limitación radica en la imposibilidad de realizar una validación externa de los modelos predictivos. Esta limitación no deriva de una omisión en el diseño, sino de la escasez actual de cohortes de replicación con características análogas (pacientes pediátricos con DT1). La validación en cohortes de adultos fue descartada deliberadamente, puesto que las profundas diferencias en la composición y dinámica del microbioma intestinal entre ambas etapas de la vida introducirían un sesgo metodológico insalvable, comprometiendo la validez de cualquier comparación.

El tercer estudio \cite{paper3}, aunque se beneficia de la inclusión de múltiples cohortes, debe considerar con cautela la generalización de sus conclusiones debido a la limitada diversidad geográfica (principalmente América del Norte y China). Además, la cohorte de descubrimiento se procesó con OTUs, mientras que las cohortes de validación utilizaron ASVs. Sin embargo, los subtipos definidos se validaron consistentemente en las cohortes externas, lo que sugiere que el impacto práctico de esta discrepancia es limitado. Aun así, el uso uniforme de ASVs podría haber identificado señales adicionales que complementaran y enriquecieran los subtipos propuestos. Finalmente, para garantizar la comparabilidad entre cohortes, se restringió el análisis a 22 especies compartidas, una decisión pragmática que podría omitir información relevante fuera de ese panel.

\textbf{Limitaciones metodológicas}

Más allá de las particularidades de cada estudio, la investigación del microbioma enfrenta fuentes de variabilidad técnica que pueden impactar significativamente los análisis posteriores. Estas fuentes surgen a lo largo del flujo de trabajo: desde el diseño experimental, donde la elección de la región del gen 16S rRNA a secuenciar o los kits de extracción de ADN son determinantes; pasando por la secuenciación, donde la tecnología empleada (Illumina, PacBio, Ion Torrent) introduce sesgos específicos; y culminando en el procesamiento bioinformático. En esta última fase, la variabilidad se acentúa por la estrategia de agrupación seguida (OTUs frente a ASVs), la elección del método para realizar dicha agrupación (DADA2, QIIME2, MOTHUR) y las bases de datos de referencia utilizadas para la asignación taxonómica (Greengenes, SILVA, RDP).

\textbf{Limitaciones intrínsecas}

Finalmente, es importante abordar una limitación de naturaleza conceptual, relacionada con el diseño transversal de los estudios. Al ofrecer una «instantánea» estática del microbioma en un momento concreto, se obtiene una simplificación necesaria para el análisis, pero que inherentemente limita la observación de la dinámica temporal y las interacciones funcionales del ecosistema. Por tanto, es preciso reconocer que analizar la biología de una enfermedad desde un único punto temporal restringe la comprensión completa de sus procesos fisiopatológicos dinámicos.

\section{Cierre de la discusión}

En el futuro, es crucial avanzar hacia la integración de plataformas analíticas para unificar datos variados, combinando información del microbioma y del anfitrión para entender mejor las interacciones biológicas y sus implicaciones en la salud. Además, investigar el microbioma como barrera en la eficacia de tratamientos y analizar la resistencia antimicrobiana facilitarán el desarrollo de estrategias innovadoras.

En resumen, las principales contribuciones de la presente Tesis Doctoral se pueden sintetizar en cuatro ejes fundamentales: la capacidad del microbioma como fuente de biomarcadores, el gran potencial del uso de técnicas de ML en el análisis del microbioma, la identificación de posibles biomarcadores microbianos con potencial de aplicación clínica en el contexto de la DT1/DT2 y el SM, y el desarrollo de una nueva estrategia de estratificación del microbioma vaginal en el contexto de la VB. Estas aportaciones evidencian la utilidad práctica del enfoque propuesto, instando a continuar explorando el microbioma como una herramienta clave para la personalización de intervenciones médicas. Este compendio, en conjunto, sienta las bases para un futuro en el que las tecnologías emergentes y la integración de datos complejos configuren el panorama de la investigación biomédica, aportando beneficios tangibles a la salud pública y a la práctica clínica moderna.